 \subsection*{Notation}
 We adopt the following notational conventions throughout the paper.
 \begin{itemize}
 \item For integers $a<b$, set $[a,b]=\{a,a+1,\ldots,b\}$.
 \item Given any quantity $\alpha$ defined using the tuple $(d_1,d_2)$ or the pair of polynomials $(P_1,P_2)$, let $\alpha'$ denote the same quantity defined using the tuple $(d'_1,d'_2)=(d_0,d_1)$ or the polynomials $(P'_1,P'_2)=(P_0,P_1)$. 
 In particular, $p'_{1,j}=p_{2,d_2-j}$ and $p'_{2,j}=p_{1,j}$ when equation~\eqref{eq:edge weights} is applied to a $(d'_1,d'_2)$-bounded grading $\omega'$ on $D_{\bfa'_m}$.
 \item Equations that will be referenced globally will be assigned numbers, those that are referenced only locally (\ie within a single proof) will be assigned symbols (e.g.\ $\dagger$ or $\ddagger$). 
 In particular, symbols labeling equations will be reused but this should not lead to any confusion. 
 \end{itemize}


%%%%%%%%%%%%%%%%%%%%%%%%%%%%
\section{Maximal Dyck Paths}
\label{sec:dyck paths}

In this section we study the recursive structure present in the maximal Dyck paths $D_m$. 
To accomplish this, we note that the vectors $\bfa_m$ can be written more explicitly in terms of two-parameter Chebyshev polynomials $u_{m,k}$ ($m,k\in \ZZ$) defined recursively~by: 
\begin{equation}
 \label{eq:chebyshev}
 u_{0,k}=0,\quad
 u_{1,k}=1,\quad
 u_{m+1,k+1}=d_ku_{m,k}-u_{m-1,k-1},
\end{equation}
where $d_k$ denotes the degree of the polynomial $P_k$ in equation~\eqref{eq:reversed polynomials}. 
Then, for $m\ge 1$, we have $\bfa_m=(u_{m,1},u_{m-1,2})$. Write $\bfa'_m=(u'_{m,1},u'_{m-1,2})=(u_{m,2},u_{m-1,1})$ and set $D'_m=D_{\bfa'_m}$ for $m\ge1$. 
\begin{remark}
 \label{rem:chebyshev}
 To see the equivalence with equation~\eqref{eq:roots recursive}, one must use the identities $u_{m,k}=u_{m,k+1}$ for $m$ odd and $d_ku_{m,k}=d_{k+1}u_{m,k+1}$ for $m$ even.
\end{remark}
We record the next simple observation for future use.
\begin{lemma}
 \label{le:Dyck path inequality}
 For positive integers $d_1,d_2$ and any integers $m,k$, we have $u_{m,k\Mk +\Mk 1}u_{m\Mk -\Mk 2,k}\Mk <u_{m-1,k+1}u_{m-1,k}$.
\end{lemma}
\begin{proof}
 We work by induction on $m$, the case $m=2$ is the trivial inequality $0<1$. 
 For $m\ge3$, we have
 \begin{align*}
 u_{m,k+1}u_{m-2,k}
 &=d_ku_{m-1,k}u_{m-2,k}-u_{m-2,k-1}u_{m-2,k}\\
 &<d_ku_{m-1,k}u_{m-2,k}-u_{m-3,k-1}u_{m-1,k}\\
 &=u_{m-1,k+1}u_{m-1,k},
 \end{align*}
 where the inequality uses induction.
 The case $m\le 1$ can be handled similarly.
\end{proof}

In order to establish a recursive structure for $D_m$ we will show that the maximal Dyck paths $D_m$ and $D'_m$ are intimately related.
\begin{lemma}
 \label{le:Dyck path recursion}
 For $m\ge1$, the following hold.
 \begin{enumerate}[label=(\alph*)]
 \item\label{lemma2.3_1} The maximal Dyck path $D'_{m+1}$ can be obtained from $D_m$ via replacing each horizontal edge, together with the $\ell$ vertical edges which immediately follow it, by $d_1-\ell$ horizontal edges followed by a vertical edge.
 \item\label{lemma2.3_2} The maximal Dyck path $D_{m+1}$ can be obtained from $D'_m$ via replacing each horizontal edge, together with the $\ell$ vertical edges which immediately follow it, by $d_2-\ell$ horizontal edges followed by a vertical edge.
 \end{enumerate}
\end{lemma}
\begin{proof}
 We prove~\ref{lemma2.3_1} as~\ref{lemma2.3_2} will immediately follow by interchanging the roles of $d_1$ and $d_2$. 
 Let $D'$ denote the lattice path obtained from $D_m$ as in~\ref{lemma2.3_1}. 
 It follows from the definition that $D'$ will contain $d_1u_{m,1}-u_{m-1,2}=u_{m+1,2}$ horizontal edges and $u_{m,1}$ vertical edges. 
 We need to show that $D'$ does not cross above the main diagonal and that it is maximal with this property. 

 Write $v'_1,\ldots,v'_{u_{m,1}}$ for the vertical edges of $D'$ and for $1\le r\le u_{m,1}$ suppose $v'_r$ is immediately preceded by exactly $\ell_r$ horizontal edges of the same height.
 Suppose there exists $t$ so that $v'_t$ passes above the main diagonal, this is equivalent to the inequality $\frac{t}{\sum\limits_{r=1}^t\ell_r}>\frac{u_{m,1}}{u_{m+1,2}}$. 
 Using the equality $u_{m+1,2}=d_1u_{m,1}-u_{m-1,2}$, this may be rewritten as 
 \begin{equation}\label{dagger1}
 \tag{$\dagger$}\frac{d_1t-\sum\limits_{r=1}^t\ell_r}{t}>\frac{u_{m-1,2}}{u_{m,1}}.
 \end{equation}
 But by construction of $D'$, we see that $d_1-\ell_r$ is the number of vertical edges immediately following the $r$-th horizontal edge of $D_m$. 
 Thus, by rewriting the numerator as $d_1t-\sum\limits_{r=1}^t\ell_r=\sum\limits_{r=1}^t(d_1-\ell_r)$ in the inequality~\eqref{dagger1}, we see that the subpath of $D_m$ containing the first $t$ horizontal edges and the vertical edges immediately following these horizontal edges will cross above the main diagonal of the rectangle $[0,u_{m,1}]\times[0,u_{m-1,2}]$, a contradiction. 
 Thus $D'$ is a Dyck path, \ie it does not pass above the main diagonal.

 To see that $D'$ is maximal, suppose there exists a lattice point $(s,t)$ strictly above $D'$ which does not lie above the main diagonal. 
 Without loss of generality, we may take $s=\sum\limits_{r=1}^t\ell_r-1$ and get the inequality $\frac{t}{\sum\limits_{r=1}^t\ell_r-1}\le\frac{u_{m,1}}{u_{m+1,2}}$. 
 Using the equality $u_{m+1,2}=d_1u_{m,1}-u_{m-1,2}$, this may be rewritten as
 \begin{equation}\label{ddagger1}
 \tag{$\ddagger$}\frac{d_1t-\sum\limits_{r=1}^t\ell_r+1}{t}\le\frac{u_{m-1,2}}{u_{m,1}}.
 \end{equation}
 Now considering the same initial segment of $D_m$ as above, we see that the point $(t,d_1t-\sum\limits_{r=1}^t\ell_r+1)$ lies strictly above $D_m$, but the inequality~\eqref{ddagger1} implies this point does not lie above the main diagonal of the rectangle $[0,u_{m,1}]\times[0,u_{m-1,2}]$, contradicting the maximality of $D_m$. 
 Thus we may conclude that $D'=D'_{m+1}$ is the maximal Dyck path in the rectangle $[0,u_{m+1,2}]\times[0,u_{m,1}]$.
\end{proof}

With this we obtain the recursive structure of $D_m$. In what follows we always assume $d_1d_2\ge4$ and take $\delta_m=\begin{cases} 1 & \text{if $d_{m-1}=1$ and $m\ne3$;}\\ 0 & \text{if $d_{m-1}\ne1$ or $m=3$.}\end{cases}$
\begin{corollary}
 \label{cor:recursive dyck path structure}
 The maximal Dyck paths $D_m$, $m\ge1$, admit the following recursive structure:
 \begin{enumerate}[label=(\alph*)]
 \item\label{coro2.4_1} $D_1$ consists of a single horizontal edge;
 \item\label{coro2.4_2} $D_2$ consists of $d_2$ horizontal edges followed by a single vertical edge;
 \item\label{coro2.4_3} for $m\ge3$, the Dyck path $D_m$ can be constructed by concatenating $d_m-1-\delta_m$ copies of $D_{m-1}$ followed by a copy of $D_{m-1}$ with its first $D_{m-2-\delta_m}$ removed.
 \end{enumerate}
\end{corollary}
For the remainder of the paper we will understand the notation $D_{m-1}\setminus D_{m-2-\delta_m}$ to mean the terminal subpath of $D_{m-1}$ obtained by removing its first copy of $D_{m-2-\delta_m}$ as in Corollary~\ref{cor:recursive dyck path structure}.
\begin{remark}
 \label{rem:primed recursive structure}
 The roles of $d_1$ and $d_2$ must be interchanged when applying Corollary~\ref{cor:recursive dyck path structure} to $D'_m$.
\end{remark}
\begin{proof}
 Parts~\ref{coro2.4_1} and~\ref{coro2.4_2} are immediate from the definitions of $D_1$ and $D_2$. 
 Part~\ref{coro2.4_3} with $m=3$ follows from Lemma~\ref{le:Dyck path recursion} and part~\ref{coro2.4_2} since $D'_2$ consists of $d_1$ horizontal edges followed by a vertical edge. 

 We establish part~\ref{coro2.4_3} by induction on $m\ge4$. 
 Notice that by Remark~\ref{rem:primed recursive structure} the claimed recursive structures of $D_m$ and $D'_{m-1}$ are the same for $m\ge5$, thus we obtain the result for $D_m$ if we know the result for $D'_{m-1}$ by applying the construction from Lemma~\ref{le:Dyck path recursion}. 
 Hence it suffices to establish the claimed recursive structure for $D_4$. 

 If $d_3\ne1$, then $\delta_4=0$ and the structure of $D_4$ is immediately deduced from Lemma~\ref{le:Dyck path recursion} and part~\ref{coro2.4_3} for $D'_3$. 
 If $d_3=1$, then $D'_2$ consists of a single horizontal edge followed by a single vertical edge and $D'_3$ consists of $d_2-1=d_4-1$ copies of $D'_2$ followed by a vertical edge. 
 Applying Lemma~\ref{le:Dyck path recursion} to $D'_3$ shows that $D_4$ consists of $d_4-2$ copies of $D_3$ followed by a copy of $D_3$ with its first horizontal edge (\ie its first $D_1$) removed. 
 This establishes part~\ref{coro2.4_3} for $D_4$ and completes the proof.
\end{proof}
\begin{corollary}
 \label{cor:initial subpath}
 For $m\ge2$, if the last edge of $D_m$ is omitted, the resulting lattice path identifies with an initial subpath of the maximal Dyck path $C_m$ obtained by concatenating $d_m$ copies of $D_{m-1}$.
\end{corollary}
\begin{proof}
 We work by induction on $m$, the case $m=2$ following immediately from Corollary~\ref{cor:recursive dyck path structure}\ref{coro2.4_2}. 

 Assume $m\ge3$.
 By part \ref{coro2.4_3} of Corollary~\ref{cor:recursive dyck path structure}, in comparing $D_m$ to $C_m$ the first $d_m-1-\delta_m$ copies of $D_{m-1}$ inside $D_m$ may be ignored and the problem reduces to comparing the final $D_{m-1}\setminus D_{m-2-\delta_m}$ subpath of $D_m$ with the maximal Dyck path $D_{m-1}$.
 For $m=3$, removing the final edge of $D_{m-1}\setminus D_{m-2-\delta_m}$ produces $d_2-1$ consecutive horizontal edges which clearly identifies with an initial subpath of $D_{m-1}$.
 Assume $m\ge4$.
 There are two cases to consider.
 \begin{itemize}
 \item If $d_{m-1}\ne1$, $D_{m-1}$ consists of $d_{m-1}-1-\delta_{m-1}$ copies of $D_{m-2}$ followed by a copy of $D_{m-2}\setminus D_{m-3-\delta_{m-1}}$. 
 It follows that comparing $D_{m-1}\setminus D_{m-2}$ with $D_{m-1}$ reduces to comparing $D_{m-2}\setminus D_{m-3-\delta_{m-1}}$ with $D_{m-2}$.
 But by induction, we know that we obtain an initial subpath of $D_{m-2}$ by removing the last edge of $D_{m-2}\setminus D_{m-3-\delta_{m-1}}$.
 \item When $d_{m-1}=1$, the maximal Dyck path $D_{m-1}$ is just $D_{m-2}\setminus D_{m-3}$. 
 But $D_{m-2}$ consists of $d_m-2$ copies of $D_{m-3}$ followed by a copy of $D_{m-3}\setminus D_{m-5}$ and so $D_{m-1}$ consists of $d_m-3$ copies of $D_{m-3}$ followed by a copy of $D_{m-3}\setminus D_{m-5}$. 
 Hence comparing $D_{m-1}\setminus D_{m-3}$ with $D_{m-1}$ reduces to comparing $D_{m-3}\setminus D_{m-5}$ with $D_{m-3}$, but by induction we know that removing the last edge of $D_{m-3}\setminus D_{m-5}$ produces an initial subpath of $D_{m-3}$.
 \end{itemize}
 The two items above show that we get an initial subpath of $D_{m-1}$ by removing the last edge of $D_{m-1}\setminus D_{m-2-\delta_m}$ and thus removing the last edge of $D_m$ produces an initial subpath of $C_m$. 
\end{proof}

Let $E_m$ denote the edges of $D_m$, where $E_m=H_m\sqcup V_m$ for horizontal edges $H_m=\{h_1,\ldots,h_{u_{m,1}}\}$ and vertical edges $V_m=\{v_1,\ldots,v_{u_{m-1,2}}\}$, both labeled in the natural order along $D_m$. 
We may describe the structure of $D_m$ as follows.
\begin{lemma}[{\cite[Lemma~3.2]{rupel2}}]
 \label{le:height and depth}
 For $m\ge2$, the following hold.
 \begin{enumerate}[label=(\alph*)]
 \item There are exactly $\hgt(h_i)\coloneqq \lfloor (i-1)u_{m-1,2}/u_{m,1}\rfloor$ vertical edges of $D_m$ preceding the horizontal edge $h_i$, call this number the \emph{height} of $h_i$;
 \item There are exactly $\dpt(v_t)\coloneqq \lceil tu_{m,1}/u_{m-1,2}\rceil$ horizontal edges of $D_m$ preceding the vertical edge $v_t$, call this number the \emph{depth} of $v_t$.
 \end{enumerate}
\end{lemma}
In the natural labeling of edges, Lemma~\ref{le:height and depth} gives $h_i=i+\lfloor (i-1)u_{m-1,2}/u_{m,1}\rfloor$ for $1\le i\le u_{m,1}$, $m\ge1$ and $v_t=t+\lceil tu_{m,1}/u_{m-1,2}\rceil$ for $1\le t\le u_{m-1,2}$, $m\ge2$. 
In particular, we see that $u_{m-1,2}<u_{m,1}$ implies $D_m$ contains no consecutive vertical edges, while $u_{m,1}<u_{m-1,2}$ implies $D_m$ contains no consecutive horizontal edges. 

For the next result, recall that we work under the assumption $d_1d_2\ge4$.
\begin{corollary}
 \label{cor:edge restrictions}
 For $m\ge2$ the following hold.
 \begin{enumerate}[label=(\alph*)]
 \item $D_m$ contains at most $1+\delta_1$ vertical edges of any given depth.
 \item $D_m$ contains no consecutive horizontal edges if and only if $d_2=1$.
 \end{enumerate}
\end{corollary}
\begin{proof}
 For $D_2$, both claims are immediate from Corollary~\ref{cor:recursive dyck path structure}\ref{coro2.4_2}.
 There are two possibilities for $D_3$.
 If $d_2=1$, the result for $D_2$ together with Corollary~\ref{cor:recursive dyck path structure}\ref{coro2.4_3} shows $D_3$ contains no consecutive horizontal edges and that the vertical edges of $D_3$ all have different depths except $v_{d_1-1}$ and $v_{d_1}$, which both have depth $d_1-1$.
 
 For $d_2>1$, the result for $D_2$ together with Corollary~\ref{cor:recursive dyck path structure}\ref{coro2.4_3} shows all vertical edges of $D_3$ have different depths.
 To see that $d_2>1$ implies there are consecutive horizontal edges in $D_3$ we need to consider two case.
 When $d_1>1$, $D_3$ begins with a copy of $D_2$ by Corollary~\ref{cor:recursive dyck path structure}\ref{coro2.4_3} and thus contains consecutive horizontal edges.
 When $d_1=1$, we must have $d_2\ge4$.
 But then $D_3$ is just $D_2\setminus D_1$ and, since $d_2\ge4$, it contains consecutive horizontal edges. 

 For $m\ge4$, both claims follow by induction using Corollary~\ref{cor:recursive dyck path structure}\ref{coro2.4_3}.
\end{proof}
Analogous statements hold for $D'_m$, with horizontal edges $H'_m=\{h'_1,\ldots,h'_{u_{m,2}}\}$ and vertical edges $V'_m=\{v'_1,\ldots,v'_{u_{m-1,1}}\}$, by interchanging the roles of $d_1$ and $d_2$.

The proof of Theorem~\ref{th:combinatorial construction} will go by induction. Towards this aim we introduce notation, following Corollary~\ref{cor:recursive dyck path structure}, which captures the recursive structure in the edges of $D_m$, $m\ge3$. 
\begin{definition}
 \label{def:subpath edges}
 For $m\ge3$ and $1\le r\le d_m-1-\delta_m$, define the following subsets of $H_m$ and $V_m$: 
 \[H_{m,r}=\{h_{(r-1)u_{m-1,1}+1},h_{(r-1)u_{m-1,1}+2},\ldots,h_{ru_{m-1,1}}\};\]
 \[V_{m,r}=\{v_{(r-1)u_{m-2,2}+1},v_{(r-1)u_{m-2,2}+2},\ldots,v_{ru_{m-2,2}}\};\]
 we identify these, for each $r$, with the horizontal and vertical edges of $D_{m-1}$. Also set 
 \[H_{m,d_m-\delta_m}=\{h_{(d_m-1-\delta_m)u_{m-1,1}+1},\ldots,h_{u_{m,1}-1},h_{u_{m,1}}\};\]
 \[V_{m,d_m-\delta_m}=\{v_{(d_m-1-\delta_m)u_{m-2,2}+1},\ldots,v_{u_{m-1,2}-1},v_{u_{m-1,2}}\};\]
 we identify these subsets with the horizontal and vertical edges of $D_{m-1}\setminus D_{m-2-\delta_m}$.

 As a notational convenience, for $1\le r\le d_m-1-\delta_m$ and $1\le i\le u_{m-1,1}$ we write $h_{i,r}\coloneqq h_{(r-1)u_{m-1,1}+i}$ and similarly $v_{t,r}\coloneqq v_{(r-1)u_{m-2,2}+t}$ for $1\le t\le u_{m-2,2}$. 
 For $u_{m-2-\delta_m,1}+1\le i\le u_{m-1,1}$, set $h_{i,d_m-\delta_m}\coloneqq h_{(d_m-1-\delta_m)u_{m-1,1}+i-u_{m-2-\delta_m,1}}$ and set $v_{t,d_m-\delta_m}\coloneqq v_{(d_m-1-\delta_m)u_{m-2,2}+t-u_{m-3-\delta_m,2}}$ for $u_{m-3-\delta_m,2}+1\le t\le u_{m-2,2}$.
\end{definition}


%%%%%%%%%%%%%%%%%%%%%%%%%%%%%%%%%%%%%%%%%%%
\section{Combinatorics of compatible pairs}
\label{sec:compatible pairs}

\noindent
Let $\omega:E_m\to\ZZ_{\ge0}$ be a $(d_1,d_2)$-bounded grading of $D_m$, $m\ge1$. 
It will be convenient to write $\omega_H$ and $\omega_V$ for the restrictions of $\omega$ to $H_m$ and to $V_m$ respectively. 
In the absence of a total grading $\omega$, we refer to the maps $\omega_H:H_m\to[0,d_1]$ and $\omega_V:V_m\to[0,d_2]$ respectively as \emph{horizontal gradings} and \emph{vertical gradings} of $D_m$. 
We will often consider $\omega$ to be the pair $(\omega_H,\omega_V)$ and refer to $\omega_H$ and $\omega_V$ as being \emph{compatible} if Definition~\ref{def:compatibility} is satisfied for $\omega$. 
Since the first condition~\eqref{eq:hgc} of Definition~\ref{def:compatibility} only involves $\omega_H$, we refer to it as the \emph{horizontal grading condition}. 
Similarly, we refer to the second condition~\eqref{eq:vgc} as the \emph{vertical grading condition}.

Write $\supp(\omega)\coloneqq \{e\in E_m:\omega(e)\ne0\}$ and call this the \emph{support} of $\omega$.
Set $\supp(\omega_H)=\supp(\omega)\cap H$ and $\supp(\omega_V)=\supp(\omega)\cap V$.
Define $|\omega|_H\coloneqq \sum\limits_{h\in H_m} \omega_H(h)$ and $|\omega|_V\coloneqq \sum\limits_{v\in V_m} \omega_V(v)$. 

\subsection{Shadow Statistics}
\label{sec:shadows}
To begin we introduce notation to gain a more delicate grasp of the compatibility conditions~\eqref{eq:hgc} and~\eqref{eq:vgc} from Definition~\ref{def:compatibility}.
For a horizontal grading $\omega_H:H_m\to[0,d_1]$ and any subpath $ee'\subset D_m$, define the \emph{horizontal shadow statistic}
\[f_{\omega_H}(ee')\coloneqq -|ee'|_V+\sum\limits_{h\in(ee')_H}\omega_H(h).\]
We also define the \emph{vertical shadow statistic}
\[f_{\omega_V}(ee')\coloneqq -|ee'|_H+\sum\limits_{v\in(ee')_V}\omega_V(v)\]
for each vertical grading $\omega_V:V_m\to[0,d_2]$. 
It immediately follows from the definitions that the shadow statistics satisfy the following additivity property with respect to concatenation of paths:
\begin{equation}
 \label{eq:shadow statistic concatenation}
 f_{\omega_H}(e_1e_3)=f_{\omega_H}(e_1e_2)+f_{\omega_H}(\overline{e}_2e_3)\quad\text{and}\quad f_{\omega_V}(e_1e_3)=f_{\omega_V}(e_1e_2)+f_{\omega_V}(\overline{e}_2e_3)
\end{equation}
for edges $e_i\in E_m$ with $e_2\in e_1e_3$.

The shadow statistics give the following alternative check for compatibility.
\begin{lemma}
 \label{le:compatibility inequality}
 Let $\omega:E_m\to\ZZ_{\ge0}$ be a compatible grading of $D_m$. For $h\in H_m$ and $v\in V_m$, the following hold:
 \begin{enumerate}[label=(\alph*)]
 \item\label{lemma3.1_1} if $f_{\omega_H}(hv)<0$, then the horizontal grading condition~\eqref{eq:hgc} is satisfied for the path $hv$;
 \item\label{lemma3.1_2} if $f_{\omega_V}(hv)<0$, then the vertical grading condition~\eqref{eq:vgc} is satisfied for the path $hv$.
 \end{enumerate}
\end{lemma}
\begin{proof}
 We prove~\ref{lemma3.1_1}, the proof of~\ref{lemma3.1_2} is similar.
 
 There is nothing to show when $\omega_H(h)=0$, so assume $h\in\supp(\omega_H)$.
 Then $f_{\omega_H}(hh)>0$ and as $e$ ranges from $h$ to $v$ the value of $f_{\omega_H}(he)$ either increases, stays the same, or decreases by 1 with each step.
 Since $f_{\omega_H}(hv)<0$, we see that $f_{\omega_H}(he)$ must eventually take the value 0 with $e\ne v$, \ie the horizontal grading condition is satisfied for the path $hv$.
\end{proof}

Apart from their relationship to the compatibility conditions~\eqref{eq:hgc} and~\eqref{eq:vgc}, the shadow statistics $f_{\omega_H}$ and $f_{\omega_V}$ encode the following important information. 
For each subpath $ee'\subset D_m$, we obtain a factor $Y_{ee'}(\omega_H,\omega_V)$ of the monomial $Y_{D_m}(\omega_H,\omega_V)$ appearing in equation~\eqref{eq:total path weights} by only multiplying the weights of edges along the path $ee'$.
\begin{lemma}
 \label{le:shadows and degrees}
 The quantities $f_{\omega_H}(ee')$ and $f_{\omega_V}(ee')$ record the total $Y$-degree and the total $X$-degree respectively of the monomial $Y_{ee'}(\omega_H,\omega_V)$.
\end{lemma}

\begin{proof}
 \looseness-1
 A horizontal edge $h \in ( e e ' )_{ H}$ contributes a factor of $p _{ 1 , \omega_{ H}  ( h )}  Y^{ \omega_{ H}  ( h ) }X^{ - 1}$  to $Y_{ e e '} ( \omega_{ H} , \omega_{ V} )$ while a vertical edge $v\in(ee')_V$ contributes a factor of\break $p_{2,d_2-\omega_V(v)}X^{\omega_V(v)+1}Y^{-1}X^{-1}$. 
 The result now follows by comparing the total $Y$- and $X$-degrees of $Y_{ee'}(\omega_H,\omega_V)$ with the definitions of $f_{\omega_H}(ee')$ and $f_{\omega_V}(ee')$ respectively.
\end{proof}

For a horizontal grading $\omega_H:H_m\to[0,d_1]$, define the \emph{local shadow path} $D(h;\omega_H)$ of a horizontal edge $h\in H_m$ to be the shortest nonempty subpath $he\subset D_m$ such that $f_{\omega_H}(he)=0$, if there is no such subpath we set $D(h;\omega_H)=hv_{u_{m-1,2}}$. 
Write $D_H(h;\omega_H)\coloneqq D(h;\omega_H)\cap H_m$ and $D_V(h;\omega_H)\coloneqq D(h;\omega_H)\cap V_m$ for the \emph{local horizontal shadow} and \emph{local vertical shadow} of $h$ with respect to $\omega_H$. 
The local shadow path $D(v;\omega_V)$ is defined similarly for a vertical edge $v\in V_m$ and a vertical grading $\omega_V:V_m\to[0,d_2]$, where $D(v;\omega_V)=h_1v$ if there is no edge $e\le v$ for which $f_{\omega_V}(ev)=0$.
The local shadows $D_H(v;\omega_V)$, $D_V(v;\omega_V)$ are defined as above. 

By definition we have $f_{\omega_H}\big(D(h;\omega_H)\big)=0$ whenever the final edge of $D(h;\omega_H)$ is not $v_{u_{m-1,2}}$. 
More importantly, writing $D(h;\omega_H)=he$, Lemma~\ref{le:compatibility inequality} together with equation~\eqref{eq:shadow statistic concatenation} imply that $f_{\omega_H}(he')>0$ and $f_{\omega_H}(e'e)<0$ for any proper subpaths $he',e'e\subset D(h;\omega_H)$. 
Thus we see for $h\in\supp(\omega_H)$ and $v\in D_V(h;\omega_H)$ that the condition~\eqref{eq:hgc} is not satisfied for the path $hv$, however for any $\omega_V$ compatible with $\omega_H$ the condition~\eqref{eq:vgc} is satisfied for $h$ and $v$. 
In particular, when $\omega_V$ is compatible with $\omega_H$, $D(v;\omega_V)$ is a proper subpath of $D(h;\omega_H)$ for any $v\in D_V(h;\omega_H)$.

Similar statements hold using the vertical shadow statistic $f_{\omega_V}$.

\subsection{Recursions}
\label{sec:recursions}
We introduce in this section a recursive construction of gradings analogous to the recursive operations on Dyck paths from Lemma~\ref{le:Dyck path recursion}. 

The \emph{shadow} of a horizontal grading $\omega_H:H_m\to[0,d_1]$ is the collection of vertical edges in the local vertical shadows of all horizontal edges, \ie $\sh(\omega_H)=\bigcup_{h\in H_m} D_V(h;\omega_H)$. 
The \emph{remote shadow} of a horizontal grading $\omega_H:H_m\to[0,d_1]$ is the subset $\rsh(\omega_H)\subset\sh(\omega_H)$ obtained by excluding for each $d$ the (up to) $\omega_H(h_d)$ vertical edges of depth $d$ immediately following $h_d$. 
The shadow and remote shadow of a vertical grading $\omega_V:V_m\to[0,d_2]$ are defined similarly.
\begin{remark}
 \label{rem:remote shadows}
 The remote shadow $\rsh(\omega_H)\subset\sh(\omega_H)$ of a horizontal grading $\omega_H$ can be described as the subset consisting of those vertical edges $v\in\sh(\omega_H)$ for which there exists a vertical grading $\omega_V$ compatible with $\omega_H$ such that $\omega_V(v)>0$.
 In particular, any vertical grading $\omega_V$ compatible with $\omega_H$ must satisfy $\omega_V(v)=0$ for $v\in\sh(\omega_H)\setminus\rsh(\omega_H)$.
\end{remark}

In order to give a relationship between gradings of $D_m$ and gradings of $D'_{m+1}$, we need to partition the remote shadows according to which local shadow contains a given edge. 
\begin{definition}
 \label{def:remote_shadows}
 Let $\omega:E_m\to\ZZ_{\ge0}$ be a grading of $D_m$.
 \begin{enumerate}[label=(\alph*)]
 \item For $1\le j\le d\le u_{m,1}$, denote by $\rsh(\omega_H)_{j;d}$ the set of $v\in\rsh(\omega_H)$ of depth $d$ such that $v\in D_V(h_j;\omega_H)$ and $h_j$ is the first horizontal edge before $v$ with this property. 
 Define the \emph{local remote shadow} of the edge $h_j$ as $\rsh(h_j;\omega_H)\coloneqq \coprod\limits_{d\in[j+1,u_{m,1}]}\rsh(\omega_H)_{j;d}$.
 \item For $0\le\ell<t\le u_{m-1,2}$, denote by $\rsh(\omega_V)_{t;\ell}$ the set of $h\in\rsh(\omega_V)$ of height $\ell$ such that $h\in D_H(v_t;\omega_V)$ and $v_t$ is the first vertical edge after $h$ with this property. 
 Define the \emph{local remote shadow} of the edge $v_t$ as $\rsh(v_t;\omega_V)\coloneqq \coprod\limits_{\ell\in[0,t-2]}\rsh(\omega_V)_{t;\ell}$.
 \end{enumerate}
\end{definition}
\begin{remark}
 By the definition of the remote shadows, it is impossible to have $d=j$ or $\ell=t-1$ in Definition~\ref{def:remote_shadows}.
\end{remark}

Lemma~\ref{le:Dyck path recursion} establishes a canonical order preserving bijection between the vertical edges $V'_{m+1}$ of $D'_{m+1}$ and the horizontal edges $H_m$ of $D_m$ which we write as $\varphi=\varphi_m:V'_{m+1}\to H_m$, $\varphi(v'_i)=h_i$ for $1\le i\le u_{m,1}$. 
Thus we obtain a bijection from $d_1$-bounded horizontal gradings of $D_m$ to $d_1$-bounded vertical gradings of $D'_{m+1}$ taking a horizontal grading $\omega_H:H_m\to[0,d_1]$ to the vertical grading $\varphi^*\omega_H:V'_{m+1}\to[0,d_1]$ given by $\varphi^*\omega_H(v'_i)=d_1-\omega_H(h_i)$.
\begin{remark}
 We will abuse notation slightly and also write $\varphi^*_m$ for the bijection between horizontal gradings of $D'_m$ and vertical gradings of $D_{m+1}$ where the roles of $d_1$ and $d_2$ need to be interchanged in the definitions above, however this abuse should not lead to any confusion.
\end{remark}

The next result shows that the remote shadows for $\omega_H$ and $\varphi^*\omega_H$ are intimately related.
\begin{lemma}[{\cite[Corollary~4.18]{rupel2}}]
 \label{le:remote shadow cardinalities}
 Let $\omega_H:H_m\to[0,d_1]$ be a horizontal grading of $D_m$.
 For $1\le j<d\le u_{m,1}$, we have $|\rsh(\omega_H)_{j;d}|=|\rsh(\varphi^*\omega_H)_{d;j-1}|$.
\end{lemma}

Thus for $1\le j<d\le u_{m,1}$ we may define a bijection $\theta_{j;d}:\rsh(\omega_H)_{j;d}\to\rsh(\varphi^*\omega_H)_{d;j-1}$ which preserves the natural order determined by distance from $h_j$ and from $v'_d$ respectively. More explicitly, as the vertical edges of $\rsh(\omega_H)_{j;d}$ are read from bottom to top the corresponding horizontal edges of $\rsh(\varphi^*\omega_H)_{d;j-1}$ are read from right to left.

For a horizontal grading $\omega_H:H_m\to[0,d_1]$, write $\cG(\omega_H)$ for the collection of all $(d_1,d_2)$-bounded gradings $\omega:E_m\to\ZZ_{\ge0}$ such that the restriction $\omega|_{H_m}$ is precisely $\omega_H$ and denote by $\cC(\omega_H)\subset\cG(\omega_H)$ the subset of compatible gradings.
Let $\cG_{\rsh}(\omega_H)\subset\cG(\omega_H)$ denote those gradings $\omega$ for which $\omega(v)=0$ whenever $v\in V_m\setminus\rsh(\omega_H)$ and write $\cC_{\rsh}(\omega_H)\coloneqq \cG_{\rsh}(\omega_H)\cap\cC(\omega_H)$.
Define analogous collections of gradings associated to a vertical grading $\omega_V:V_m\to[0,d_2]$.

Define a map $\Omega=\Omega_m:\cG_{\rsh}(\omega_H)\to\cG_{\rsh}(\varphi^*\omega_H)$ as follows:
\[\Omega(\omega_V)(h')=
 \begin{cases}
 0 & \text{ if $h'\in H'_{m+1}\setminus\rsh(\varphi^*\omega_H)$;}\\
 \omega_V(v) & \text{ if $h'=\theta_{j;d}(v)$ for $v\in\rsh(\omega_H)_{j;d}$.}
 \end{cases}\]
Note that $\Omega$ admits an obvious inverse map. 
\begin{remark}
 Given a grading $\omega:E_m\to\ZZ_{\ge0}$ of $D_m$ where $\omega_V\notin\cG_{\rsh}(\omega_H)$, the map $\Omega$ may still be applied to $\omega_V$ to produce a horizontal grading in $\cG_{\rsh}(\varphi^*\omega_H)$. 
 This observation will be used without mention in the statements of Lemma~\ref{le:f_and_Omega} and Proposition~\ref{prop:piecewise equivalence} as well as in the proof of Corollary~\ref{cor:support and remote shadow}.
\end{remark}

The following result shows that we have some control over the shadow statistics under this operation.
It is also the essential ingredient for understanding the piecewise compatible gradings introduced in the next section.
\begin{lemma}[{\cite[Lemma~4.19]{rupel2}}]
 \label{le:f_and_Omega}
 Let $\omega:E_m\to\ZZ_{\ge0}$ be a grading on $D_m$.
 Suppose $h'=\theta_{j;d}(v)$ for a vertical edge $v\in\rsh(\omega_H)_{j;d}\cap\supp(\omega_V)$. Then $f_{\Omega(\omega_V)}(h'v'_d)=f_{\omega_V}(h_jv)$.
\end{lemma}
This crucial result also shows that $\Omega$ restricts to a map $\cC_{\rsh}(\omega_H)\to\cC_{\rsh}(\varphi^*\omega_H)$, \ie that the pair $(\Omega(\omega_V),\varphi^*\omega_H)$ gives a compatible grading of $D'_{m+1}$ exactly when $\omega_V\in\cC_{\rsh}(\omega_H)$.
\begin{proposition}[{\cite[Lemma~4.20]{rupel2}}]
 \label{prop:compatibility equivalence}
 Let $\omega_H:H_m\to[0,d_1]$ be a horizontal grading of $D_m$. 
 For a vertical grading $\omega_V\in\cG_{\rsh}(\omega_H)$, we have $\omega_V\in\cC_{\rsh}(\omega_H)$ if and only if $\Omega(\omega_V)\in\cC_{\rsh}(\varphi^*\omega_H)$.
\end{proposition}

\subsection{Piecewise Compatibility}
\label{sec:piecewise compatibility}
Our goal in this section is to understand which gradings on $D_m$, $m\ge3$, are obtained by gluing together compatible gradings on the $D_{m-1}$ subpaths of $D_m$ found in Corollary~\ref{cor:recursive dyck path structure}\ref {coro2.4_3}.
\begin{definition}
 \label{def:agregate grading}
 Fix $m\ge3$. Consider $(d_1,d_2)$-bounded compatible gradings $\omega_r=(\omega_{H,r},\omega_{V,r})$ of $D_{m-1}$ for $1\le r\le d_m-\delta_m$. We assume
 \begin{equation}
 \label{eq:V restriction}
 \omega_{V,d_m-\delta_m}(v)=0
 \end{equation}
 for $v$ in the first $D_{m-2-\delta_m}$ subpath of $D_{m-1}$ and
 \begin{equation}
 \label{eq:H restriction}
 \omega_{H,d_m-\delta_m}(h)=\ell
 \end{equation}
 for $h$ in the first $D_{m-2-\delta_m}$ subpath of $D_{m-1}$ if $h$ is immediately followed by exactly $\ell$ vertical edges inside $D_{m-2-\delta_m}$. 

 Define a grading $\omega:E_m\to\ZZ_{\ge0}$ of $D_m$ by
 \[\omega(e)=\begin{cases}\omega_{H,r}(e) & \text{if $e\in H_{m,r}$;}\\ \omega_{V,r}(e) & \text{if $e\in V_{m,r}$;}\end{cases}\]
 where we identify subsets of edges in $D_m$ with edges of $D_{m-1}$ as in Definition~\ref{def:subpath edges}. We will refer to any grading on $D_m$ obtained in this way as \emph{piecewise compatible}.
\end{definition}
\begin{remark}
 \label{rem:restricted gradings}
 Every compatible grading of $D_m$, $m\ge3$, is piecewise compatible. 
 Given any grading $\omega$ of $D_m$ and $1\le r\le d_m-\delta_m$, we will denote by $\omega_r=(\omega_{H,r},\omega_{V,r})$ the grading of $D_{m-1}$ obtained by restricting $\omega$ to the $r$-th copy of $D_{m-1}$ inside $D_m$, where $\omega_{d_m-\delta_m}=(\omega_{H,d_m-\delta_m},\omega_{V,d_m-\delta_m})$ denotes the grading on $D_{m-1}$ satisfying the conditions~\eqref{eq:V restriction} and~\eqref{eq:H restriction} of Definition~\ref{def:agregate grading}.
\end{remark}
\begin{remark}
 \label{rem:alternate restrictions}
 When considering piecewise compatible gradings $\omega:E'_m\to\ZZ_{\ge0}$ of $D'_m$ we will instead make the following assumptions on the gradings $\omega_{H',d'_m-\delta'_m}$ and $\omega_{V',d'_m-\delta'_m}$ of $D'_{m-1}$:
 \[\omega_{H',d'_m-\delta'_m}(h')=0\]
 for $h'$ in the first $D'_{m-2-\delta'_m}$ subpath of $D'_{m-1}$ and
 \[\omega_{V',d'_m-\delta'_m}(v')=d\]
 for $v'$ in the first $D'_{m-2-\delta'_m}$ subpath of $D'_{m-1}$ if $v'$ is immediately preceded by exactly $d$ horizontal edges inside $D'_{m-2-\delta'_m}$. 
\end{remark}

The next result shows that only the final edge of $D_m$ needs to be considered in order to verify (global) compatibility of a piecewise compatible grading.
\begin{lemma}
 \label{le:compatibility check}
 Let $\omega:E_m\to\ZZ_{\ge0}$ be a piecewise compatible grading on $D_m$, $m\ge3$.
 Then one of the compatibility conditions~\eqref{eq:hgc} or~\eqref{eq:vgc} is satisfied for every $h\in H_m$ and every $v\in V_m\setminus\{v_{u_{m-1,2}}\}$.
 In particular, a piecewise compatible grading on $D_m$ is compatible if and only if one of the compatibility conditions~\eqref{eq:hgc} or~\eqref{eq:vgc} holds for all paths $hv_{u_{m-1,2}}$ with $h\in\bigsqcup\limits_{r=1}^{d_m-1-\delta_m} H_{m,r}$.
\end{lemma}
\begin{proof}
 Following~\cite[Remark~2.22]{rupel2}, we have a principle of non-interaction between adjacent $D_{m-1}$ subpaths of $D_m$. 
 More precisely, one of the compatibility conditions~\eqref{eq:hgc} or~\eqref{eq:vgc} will always be satisfied for paths $hv$ with $h\in H_{m,r}$ and $v\in V_{m,s}$ for $1\le r<s\le d_m-1-\delta_m$. 
 Since each pair $(\omega_{H,r},\omega_{V,r})$ is compatible, it only remains to verify a compatibility condition for $h\in\bigsqcup\limits_{r=1}^{d_m-1-\delta_m} H_{m,r}$ and $v\in V_{m,d_m-\delta_m}$. 

 By Corollary~\ref{cor:initial subpath}, we may again apply~\cite[Remark~2.22]{rupel2} to see that one of the compatibility conditions will always be satisfied for all $v\in V_{m,d_m-\delta_m}$ with $v\ne v_{u_{m-1,2}}$.
 Thus a compatibility conditions only needs to be verified for paths $hv_{u_{m-1,2}}$ with $h\in\bigsqcup\limits_{r=1}^{d_m-1-\delta_m} H_{m,r}$ to verify compatibility of $\omega$.
\end{proof}
\begin{corollary}
 \label{cor:automatic compatibility}
 When $d_m=1$, every piecewise compatible grading of $D_m$, $m\ge3$, is compatible.
\end{corollary}
\begin{proof}
 When $d_m=1$, the set of horizontal edges $\bigsqcup\limits_{r=1}^{d_m-1-\delta_m} H_{m,r}$ is empty and thus the compatibility condition of Lemma~\ref{le:compatibility check} is trivially satisfied.
\end{proof}

Next we observe that piecewise compatible gradings are well-behaved under the operations $\varphi^*$ and $\Omega$ introduced in Section~\ref{sec:recursions}.
\begin{proposition}
 \label{prop:piecewise equivalence}
 Let $\omega:E_m\to\ZZ_{\ge0}$ be a $(d_1,d_2)$-bounded grading on $D_m$ for $m\ge3$. 
 Then $\omega$ is piecewise compatible if and only if $(\Omega(\omega_V),\varphi^*\omega_H)$ is piecewise compatible.
\end{proposition} 
\begin{proof}
 We prove the forward implication, the other direction can be obtained by reversing the argument. 

 Assume $\omega$ is piecewise compatible and, for $1\le r\le d_m-\delta_m$, consider $h'\in H'_{m+1,r}$ and $v'_t\in V'_{m+1,r}$ with $h'<v'_t$. 
 If $h'\notin D(v'_t;\varphi^*\omega_H)$, then the vertical grading condition~\eqref{eq:vgc} is satisfied for the path $h'v'_t$. 
 So we assume $h'\in D(v'_t;\varphi^*\omega_H)$ and need to show that the horizontal grading condition~\eqref{eq:hgc} is satisfied for the path $h'v'_t$.

 For $\Omega(\omega_V)(h')=0$, there is nothing to check so assume $h'\in\supp(\Omega(\omega_V))$.
 Set $j-1=\hgt(h')$ so that $h'\in\rsh(\varphi^*\omega_H)_{d;j-1}$ with $j<d\le t$.
 Then $\Omega(\omega_V)(h')=\omega_V(v)$, where $v\in\rsh(\omega_H)_{j;d}\cap\supp(\omega_V)$ with $h'=\theta_{j;d}(v)$.

 Note that $h_j\in H_{m,r}$ and $v\in V_{m,r}$, thus by piecewise compatibility the vertical grading condition~\eqref{eq:vgc} is satisfied for the path $h_jv$. 
 That is, there exists $e\in\overline{h}_jv$ so that $f_{\omega_V}(ev)=0$.
 By piecewise compatibility, each vertical edge in $h_j\overline{e}$ also satisfies the vertical grading condition with $h_j$.
 It follows that $f_{\omega_V}(h_jv)<0$.

 By Lemma~\ref{le:f_and_Omega}, we thus have $f_{\Omega(\omega_V)}(h'v'_d)<0$ and so the horizontal grading condition is satisfied for the path $h'v'_d$ by Lemma~\ref{le:compatibility inequality}.
 Since $h'v'_d$ is an initial subpath of $h'v'_t$, the horizontal grading condition is also satisfied for the path $h'v'_t$. 
 Since $h'$ and $v'_t$ were arbitrary, we see that $(\Omega(\omega_V),\varphi^*\omega_H)$ is piecewise compatible.
\end{proof} 

We aim now to understand precisely when compatibility fails for a piecewise compatible grading. 
The definition below provides the necessary conditions for a piecewise compatible grading $\omega$ constructed as in Definition~\ref{def:agregate grading} to be incompatible.
\begin{definition}
 \label{def:blocking and justification}
 Let $\omega_H:H_m\to[0,d_1]$ be a horizontal grading on $D_m$, $m\ge3$. We say a horizontal edge $h\in H_m$ is \emph{blocking for $\omega_H$} if the following hold:
 \begin{itemize}
 \item $h\in H_m\setminus H_{m,d_m-\delta_m}$;
 \item $D(h;\omega_H)=hv_{u_{m-1,2}}$;
 \item $h$ is the maximal (\ie furthest to the right) horizontal edge with these properties.
 \end{itemize}
 We call $\omega_H$ \emph{left-justified} at a blocking edge $h_i\in H_m$ if there exists $k\ge i$ so that $\omega_H(h_j)>0$ for $i\le j\le k$ and $\omega_H(h_j)=0$ for $j>k$. 
 Such a horizontal grading is \emph{strongly left-justified at $h_i$} if in addition the following hold:
 \begin{itemize}
 \item $\omega_H(h_j)=d_1$ for $i\le j<k$;
 \item $f_{\omega_H}(h_iv_{u_{m-1,2}})=0$.
 \end{itemize}
 Let $\omega_V:V_m\to[0,d_2]$ be a vertical grading on $D_m$, $m\ge3$. 
 For a horizontal edge $h_i\in H_m$, $\omega_V$ is called \emph{right-justified with respect to $h_i$} if there is a vertical edge $v_s\in h_iv_{u_{m-1,2}}$ so that $\omega_V(v_t)>0$ for $s\le t\le u_{m-1,2}$ and $\omega_V(v_t)=0$ for all vertical edges $v_t\in(h_i\overline{v}_s)_V$. 
 Such a vertical grading is \emph{strongly right-justified with respect to $h_i$} if in addition the following hold:
 \begin{itemize}
 \item $\omega_V(v_t)=d_2$ for $s<t\le u_{m-1,2}$;
 \item $D(v_{u_{m-1,2}};\omega_V)=h_iv_{u_{m-1,2}}$ with $f_{\omega_V}(h_iv_{u_{m-1,2}})=0$.
 \end{itemize}
\end{definition}

\begin{proposition}
 \label{prop:strong justification implication}
 Let $\omega:E_m\to\ZZ_{\ge0}$ be a piecewise compatible grading of $D_m$, $m\ge3$, for which $\omega_H$ admits the blocking edge $h_i\in H_m$ and $\omega_V$ is strongly right-justified with respect to $h_i$.
 Then $\omega_H$ is left-justified at $h_i$ and $\supp(\omega_H)\cap h_iv_{u_{m-1,2}}=\rsh(\omega_V)\cap h_iv_{u_{m-1,2}}$.
\end{proposition}
\begin{proof}
 Since $\omega_V$ is strongly right-justified with respect to $h_i$, we have $D(v_{u_{m-1,2}};\omega_V)=h_iv_{u_{m-1,2}}$ and thus $\supp(\omega_H)\cap h_iv_{u_{m-1,2}}\subset\rsh(\omega_V)\cap h_iv_{u_{m-1,2}}$.
 To see equality of these sets we show that they must have the same cardinality.

 As $\omega_V$ is strongly right-justified with respect to $h_i$, we have $\omega_V(v)=d_2$ for each vertical edge $v\in\overline{v}_sv_{u_{m-1,2}}$, where $s=u_{m-1,2}-\left\lfloor\frac{u_{m,1}-i+1}{d_2}\right\rfloor$.
 This implies $\omega_H(h)=0$ for each horizontal edge $h\in\overline{v}_sv_{u_{m-1,2}}$.
 Otherwise both grading conditions would fail for the path $hv$, where $v$ is the first vertical edge after $h$, a contradiction with piecewise compatibility of $\omega$.

 Then observe that the vertical edge $v_s$ has depth $\left\lceil\frac{\left(u_{m-1,2}-\left\lfloor\frac{u_{m,1}-i+1}{d_2}\right\rfloor\right)u_{m,1}}{u_{m-1,2}}\right\rceil$ by Lemma~\ref{le:height and depth}.
 Using once more that $\omega_V$ is strongly right-justified with respect to $h_i$, we must also have $\omega_H(h)=0$ for each of the $u_{m,1}-i+1-d_2\left\lfloor\frac{u_{m,1}-i+1}{d_2}\right\rfloor$ horizontal edges $h$ immediately preceding $v_s$ in order for piecewise compatibility to hold.

 The above discussion has shown that $\omega_H(h_j)=0$ whenever $j$ is larger than the following quantity:
 \begin{align*}
 &\left\lceil\frac{\left(u_{m-1,2}-\left\lfloor\frac{u_{m,1}-i+1}{d_2}\right\rfloor\right)u_{m,1}}{u_{m-1,2}}\right\rceil-\left(u_{m,1}-i+1-d_2\left\lfloor\frac{u_{m,1}-i+1}{d_2}\right\rfloor\right)\\
 &\quad=i-1+\left\lceil\frac{-\left\lfloor\frac{u_{m,1}-i+1}{d_2}\right\rfloor u_{m,1}}{u_{m-1,2}}\right\rceil+d_2\left\lfloor\frac{u_{m,1}-i+1}{d_2}\right\rfloor\\
 &\quad=i-1+\left\lceil\frac{\left\lfloor\frac{u_{m,1}-i+1}{d_2}\right\rfloor u_{m-2,1}}{u_{m-1,2}}\right\rceil,
 \end{align*}
 where both equalities follow from the identity $\lceil n+x\rceil=n+\lceil x\rceil$ which holds for all real numbers $x$ and all integers $n$.
 This discussion also shows that $h_iv_{u_{m-1,2}}\cap\big(\sh(\omega_V)\setminus\rsh(\omega_V)\big)=(h_{i+d}v_{u_{m-1,2}})_H$, where $d=\left\lceil\frac{\left\lfloor\frac{u_{m,1}-i+1}{d_2}\right\rfloor u_{m-2,1}}{u_{m-1,2}}\right\rceil$. 
 Since $D(v_{u_{m-1,2}};\omega_V)=h_iv_{u_{m-1,2}}$, it follows that 
 \[|\rsh(\omega_V)\cap h_iv_{u_{m-1,2}}|=\left\lceil\frac{\left\lfloor\frac{u_{m,1}-i+1}{d_2}\right\rfloor u_{m-2,1}}{u_{m-1,2}}\right\rceil.\]

 Now observe the inequality 
 \[\left\lceil\frac{\left\lfloor\frac{u_{m,1}-i+1}{d_2}\right\rfloor u_{m-2,1}}{u_{m-1,2}}\right\rceil\le\left\lceil\frac{(u_{m,1}-i+1)u_{m-2,1}}{d_2u_{m-1,2}}\right\rceil=\left\lceil\frac{(u_{m,1}-i+1)u_{m-2,2}}{d_1u_{m-1,1}}\right\rceil,\] 
 where the equality can be deduced from the identities in Remark~\ref{rem:chebyshev}.
 By Lemma~\ref{le:Dyck path inequality}, the last expression above is not larger than 
 \begin{equation}
 \label{eq:strongly justified support}
 \left\lceil\frac{(u_{m,1}-i+1)u_{m-1,2}}{d_1u_{m,1}}\right\rceil=\left\lceil\frac{u_{m-1,2}-\left\lfloor\frac{(i-1)u_{m-1,2}}{u_{m,1}}\right\rfloor}{d_1}\right\rceil,
 \end{equation}
 where the equality follows from right to left using the identities $-\lfloor x\rfloor=\lceil-x\rceil$, $\lceil n+x\rceil=n+\lceil x\rceil$, and $\left\lceil\frac{\lceil x\rceil}{n}\right\rceil=\left\lceil\frac{x}{n}\right\rceil$ which hold for all real numbers $x$ and all positive integers $n$.
 
 But $h_i$ is blocking and $\omega_H$ is $d_1$-bounded so that
 \[|\supp(\omega_H)\cap h_iv_{u_{m-1,2}}|\ge\left\lceil\frac{|h_iv_{u_{m-1,2}}|_V}{d_1}\right\rceil=\left\lceil\frac{u_{m-1,2}-\left\lfloor\frac{(i-1)u_{m-1,2}}{u_{m,1}}\right\rfloor}{d_1}\right\rceil.\]
 Combining this observation with the inequalities leading up to equation~\eqref{eq:strongly justified support}, we see that
 \begin{align*}
 |\supp(\omega_H)\cap h_iv_{u_{m-1,2}}|
 &\ge\left\lceil\frac{u_{m-1,2}-\left\lfloor\frac{(i-1)u_{m-1,2}}{u_{m,1}}\right\rfloor}{d_1}\right\rceil\\
 &\ge\left\lceil\frac{\left\lfloor\frac{u_{m,1}-i+1}{d_2}\right\rfloor u_{m-2,1}}{u_{m-1,2}}\right\rceil=|\rsh(\omega_V)\cap h_iv_{u_{m-1,2}}|.
 \end{align*}
 But either inequality being strict is impossible since $\supp(\omega_H)$ does not intersect $\sh(\omega_V)\setminus\rsh(\omega_V)$. 
 Thus $\omega_H$ must be left-justified at $h_i$ with 
 \begin{align}
 \label{eq:incompatible support condition}
 |\supp(\omega_H)\cap h_iv_{u_{m-1,2}}|
 &=\left\lceil\frac{u_{m-1,2}-\left\lfloor\frac{(i-1)u_{m-1,2}}{u_{m,1}}\right\rfloor}{d_1}\right\rceil\\
 \nonumber &=\left\lceil\frac{\left\lfloor\frac{u_{m,1}-i+1}{d_2}\right\rfloor u_{m-2,1}}{u_{m-1,2}}\right\rceil=|\rsh(\omega_V)\cap h_iv_{u_{m-1,2}}|.
 \end{align}
 In particular, $\supp(\omega_H)\cap h_iv_{u_{m-1,2}}=\rsh(\omega_V)\cap h_iv_{u_{m-1,2}}$.
\end{proof}
\begin{remark}
 The middle equality of equation~\eqref{eq:incompatible support condition} does not hold for all $i$, this equality is a consequence of the hypotheses and thus provides a necessary condition for the existence of a piecewise compatible grading as in Proposition~\ref{prop:strong justification implication}.
\end{remark} 

The next result will show that this condition is also sufficient and that such gradings are the only piecewise compatible gradings which are not compatible.
\begin{theorem}
 \label{th:blocking edge conditions}
 Let $\omega:E_m\to\ZZ_{\ge0}$ be a piecewise compatible grading of $D_m$, $m\ge3$.
 \begin{enumerate}[label=(\alph*)]
 \item\label{theo3.20_1} If $\omega_H$ does not admit a blocking edge, then $\omega$ is compatible.
 \item\label{theo3.20_2} Suppose $\omega_H$ admits a blocking edge $h_i$ but $\omega$ is not compatible.
 Then the following hold:
 \begin{enumerate}[(\roman*)]
 \item\label{theo3.20_2_a} $D(v_{u_{m-1,2}};\omega_V)=h_iv_{u_{m-1,2}}$ with $f_{\omega_V}(h_iv_{u_{m-1,2}})=0$;
 \item\label{theo3.20_2_b} $\omega_H$ is left-justified at $h_i$ and $\omega_V$ is strongly right-justified with respect to $h_i$.
 \end{enumerate}
 If in addition $m\ge4$, the following also hold:
 \begin{enumerate}[(\roman*)]\addtocounter{enumii}{2}
 \item\label{theo3.20_2_c} $f_{\omega_H}(h_iv_{u_{m-1,2}})=0$;
 \item\label{theo3.20_2_d} $\omega_H$ must be strongly left-justified at $h_i$;
 \item\label{theo3.20_2_e} $\hgt(h_{i+1})=\hgt(h_i)+\delta_1$ when $|\supp(\omega_H)\cap h_iv_{u_{m-1,2}}|>1$.
 \end{enumerate}
 \end{enumerate}
\end{theorem}
\begin{remark}
 When $d_m=1$, the hypotheses of Theorem~\ref{th:blocking edge conditions} cannot apply by Corollary~\ref{cor:automatic compatibility}.
\end{remark}
\begin{proof}
 If $\omega_H$ does not admit a blocking edge, any horizontal edge $h\in H_m$ has a local shadow path of the form $D(h;\omega_H)=he$ with $e<v_{u_{m-1,2}}$, \ie the horizontal grading condition is satisfied for $h$ and $v_{u_{m-1,2}}$. By Lemma~\ref{le:compatibility check}, this implies $\omega$ is compatible. 
 This establishes~\ref{theo3.20_1}. 
 
 From now on we assume $\omega_H$ admits a blocking edge $h_i$ and $d_m\ne1$. 
 There are two possible cases.
 First consider the case $\hgt(h_i)\ge u_{m-1,2}-d_1$.
 Since by definition $h_i$ is not contained in the final $D_{m-1}\setminus D_{m-2-\delta_m}$ subpath of $D_m$, we must have $d_1\ge2$ and so this case can occur only if one of the following holds:
 \begin{itemize}
 \item $m=3$ with $d_1\ne1$;
 \item $m=4$ with $d_1,d_2\ne1$;
 \item $m=5$ with $d_2=1$.
 \end{itemize}
 When $m=4$ or $m=5$ above, we must have $h_i=h_{u_{m-1,1},d_m-1-\delta_m}$ (see Definition~\ref{def:subpath edges} for notation).

 Let $\ell=\hgt(h_i)$.
 Then, since $d_1\ne1$ in all cases above, each vertical edge in $\overline{v}_{\ell+1}\overline{v}_{u_{m-1,2}}$ if $m=3$ (resp. in $h_{i+1}\overline{v}_{u_{m-1,2}}$ if $m=4$ or $m=5$) is immediately preceded by exactly $d_2$ horizontal edges inside $h_iv_{u_{m-1,2}}$ while $v_{u_{m-1,2}}$ is immediately preceded by exactly $d_2-1$ horizontal edges.
 In particular, we see that the horizontal grading condition fails for the path $h_iv_{u_{m-1,2}}$ exactly when:
 \begin{itemize}
 \item $\omega_V(v_{\ell+1})=\dpt(v_{\ell+1})-i$ and $\omega_V(v)=d_2$ for $v\in(\overline{v}_{\ell+1}v_{u_{m-1,2}})_V$ for $m=3$;
 \item $\omega_V(v)=d_2$ for all $v\in(h_iv_{u_{m-1,2}})_V$ for $m=4$ or $m=5$.
 \end{itemize}
 In either case we have $D(v_{u_{m-1,2}};\omega_V)=h_iv_{u_{m-1,2}}$ with $f_{\omega_V}(h_iv_{u_{m-1,2}})=0$.
 Note that in each of the cases above, $\omega_H$ is left-justified at $h_i$ with $k=i$ in Definition~\ref{def:blocking and justification} and $\omega_V$ is strongly right-justified with respect to $h_i$.
 This establishes the claims in the first part of~\ref{theo3.20_2} for these cases.
 Observe that our assumptions when $m=4$ or $m=5$ imply $f_{\omega_H}(h_iv_{u_{m-1,2}})=0$ and that $\omega_H$ is strongly left-justified at $h_i$.
 Since $\supp(\omega_H)\cap h_iv_{u_{m-1,2}}=\{h_i\}$, this establishes the second part of~\ref{theo3.20_2} in these cases.
 
 Now assume $m\ge4$ and $\hgt(h_i)<u_{m-1,2}-d_1$. Then there must exist $j>i$ so that $D(h_j;\omega_H)=h_jv_{u_{m-1,2}-\ell}$ with $1\le\ell\le\omega_H(h_i)-\delta_1$
 (the extra $\delta_1$ must be included here since $d_2=1$ implies all horizontal edges of $D_m$ have different heights, in other words $d_2=1$ implies $h_i$ is immediately followed by a vertical edge). 
 Assume that $j$ is chosen so that $\ell$ is minimal, in particular when $d_1=1$ we must have $\ell=1$.

 By Lemma~\ref{le:compatibility check}, the vertical grading condition must be satisfied for the paths $h_iv$ with $v\in(h_i\overline{v}_{u_{m-1,2}})_V$.
 For each such $v$, we have $D(v;\omega_V)=h_{j(v)}v$ for some $j(v)>i$, in particular $f_{\omega_V}(h_{j(v)}v)=0$. 
 Since $h_i$ is blocking, it cannot be contained in the shadow of any of these vertical edges.
 Moreover, when $d_1=1$, the edge $h_j$ will also not be contained in the shadow of any of these vertical edges. 
 Thus we see that there are at least $1+\delta_2$ horizontal edges of the path $h_iv_{u_{m-1,2}-1}$ lying outside the shadows of its vertical edges and applying equation~\eqref{eq:shadow statistic concatenation} shows $f_{\omega_V}(h_iv_{u_{m-1,2}-1})\le-(1+\delta_2)$.
 But by Corollary~\ref{cor:recursive dyck path structure} there are $d_2-1-\delta_2$ horizontal edges immediately preceding $v_{u_{m-1,2}}$ and, since $\omega_V(v_{u_{m-1,2}})\le d_2$, we must have $f_{\omega_V}(h_iv_{u_{m-1,2}})\le0$. 
 We conclude that one of the following holds:
 \begin{itemize}
 \item $D(v_{u_{m-1,2}};\omega_V)$ is a proper subpath of $h_iv_{u_{m-1,2}}$ by Lemma~\ref{le:compatibility inequality} and thus $\omega$ is compatible;
 \item $D(v_{u_{m-1,2}};\omega_V)=h_iv_{u_{m-1,2}}$ with $f_{\omega_V}(h_iv_{u_{m-1,2}})=0$ and both compatibility conditions fail for the path $h_iv_{u_{m-1,2}}$.
 \end{itemize}
 This establishes claim~\ref{theo3.20_2_a} of~\ref{theo3.20_2} in this case. 
 When $d_1=1$, we must have $f_{\omega_H}(h_iv_{u_{m-1,2}})=0$ for otherwise $h_i$ could not be blocking.
 This gives claim~\ref{theo3.20_2_c} of~\ref{theo3.20_2} when $d_1=1$. 
 To complete the proof of~\ref{theo3.20_2_c} for $d_1>1$ and $m\ge4$, we observe that $h_i$ being a blocking edge implies $f_{\omega_H}(h_iv_{u_{m-1,2}})\ge0$.
 Our aim then is to show that $f_{\omega_H}(h_iv_{u_{m-1,2}})>0$ implies the second situation above is impossible.

 Indeed, $f_{\omega_H}(h_iv_{u_{m-1,2}})>0$ can only occur if we take $\ell\le\omega_H(h_i)-1-\delta_1$ above.
 But, assuming $d_1>1$ and $m\ge4$, there are $d_2$ horizontal edges of $D_m$ immediately preceding each of the vertical edges 
 \[v_{u_{m-1,2}-d_1+2+\delta_1},v_{u_{m-1,2}-d_1+3+\delta_1},\ldots,v_{u_{m-1,2}-1},\]
 and $d_2-1$ horizontal edges immediately preceding $v_{u_{m-1,2}}$ (by Corollary~\ref{cor:recursive dyck path structure}, the terminal subpath of $D_m$ containing all these edges identifies with the terminal subpath $D_3\setminus D_2$ inside $D_3$). 
 It follows that $D(v_{u_{m-1,2}};\omega_V)$ must be a subpath of $h_jv_{u_{m-1,2}}$ and so the vertical grading condition is satisfied for the path $h_iv_{u_{m-1,2}}$. 
 In particular, $\omega$ is compatible by Lemma~\ref{le:compatibility check}, this completes the proof of~\ref{theo3.20_2_c}.

 The arguments above also establish the following when $m\ge4$, $d_m\ne1$, and $\hgt(h_i)<u_{m-1,2}-d_1$:
 \begin{itemize}
 \item if $\omega_H(h_i)<d_1$ or $h_i$ is immediately followed by $1+\delta_1$ vertical edges, then either $d_1=1$ and $h_i$ cannot possibly be blocking or there must exist a horizontal edge $h_j$ as in the previous paragraph and compatibility again holds, this gives~\ref{theo3.20_2_e} once we have established~\ref{theo3.20_2_d}, \ie once we know that $\omega_H$ is strongly left-justified at $h_i$; 
 \item if $\omega_V(v_{u_{m-1,2}-t})<d_2$ for any $0\le t\le d_1-\delta_1$, then the piecewise compatible grading $\omega$ must be compatible.
 \end{itemize}

 We prove~\ref{theo3.20_2_c} and~\ref{theo3.20_2_d} by induction on $m\ge3$, $d_m\ne1$.
 The base case $m=3$ of~\ref{theo3.20_2_b} was established in the first part of the proof.
 Suppose $m\ge4$ and $\omega$ is not compatible. 
 By Proposition~\ref{prop:piecewise equivalence} the grading $\big((\varphi^*_{m-1})^{-1}\omega_V,\Omega_{m-1}^{-1}(\omega_H)\big)=:(\omega_{H'},\omega_{V'})$ of $D'_{m-1}$ is piecewise compatible, but not compatible by Proposition~\ref{prop:compatibility equivalence}. 
 By part~\ref{theo3.20_1}, there must be a blocking edge $h'_j$ for $\omega_{H'}$.
 Applying~\ref{theo3.20_2_b} to the grading $(\omega_{H'},\omega_{V'})$ we see that $\omega_{H'}$ is left-justified at $h'_j$ and $\omega_{V'}$ is strongly right-justified with respect to this blocking edge.

 When $m=4$, we have $\supp(\omega_{H'})\cap h'_jv'_{u'_{m-2,2}}=\{h'_j\}$ and from the definition of $\varphi^*_{m-1}$ we see that $\omega_V$ is strongly right-justified with respect to $h_i$.
 This requires the extra observation above that we had to take $k=i$ in the definition of left-justification for the case $m=3$.
 For $m\ge5$, claim~\ref{theo3.20_2_d} applied to the grading $(\omega_{H'},\omega_{V'})$ shows that $\omega_{H'}$ is strongly left-justified at $h'_j$ and again the definition of $\varphi^*_{m-1}$ shows that $\omega_V$ is strongly right-justified with respect to $h_i$. 
 By Proposition~\ref{prop:strong justification implication}, we see that $\omega_H$ must be left-justified at $h_i$. 

 It remains to argue that $\omega_H$ is strongly left-justified at $h_i$, but this is immediate from Lemma~\ref{le:remote shadow cardinalities} and the definition of the maps $\theta$. 
 Indeed, since $\omega_{H'}$ is strongly left-justified at its blocking edge $h'_j$, the remote shadows of the horizontal edges in $h'_jv'_{u_{m-2,1}}$ are linearly ordered in the opposite order to the horizontal edges in $\supp(\omega_{H'})\cap h'_jv'_{u_{m-2,1}}$. 
 Since $\omega_V$ is strongly right-justified with respect to $h_i$, analogous statements can be made about the remote shadows of the vertical edges in $h_iv_{u_{m-1,2}}$. 
 But the maps $\theta$ are compatible with these orderings and so $\omega_{V'}$ being strongly right-justified with respect to $h'_j$ forces $\omega_H=\Omega_{m-1}(\omega_{V'})$ to be strongly left-justified at $h_i$.
 This completes the proof of~\ref{theo3.20_2_b} and~\ref{theo3.20_2_d}.
\end{proof}

The next result severely restricts which horizontal edges can be blocking.
\begin{corollary}
 Let $\omega:E_m\to\ZZ_{\ge0}$ be a piecewise compatible grading of $D_m$, $m\ge5$, which is not compatible.
 Write $h_i\in H_m$ for the blocking edge of $\omega_H$.
 Then either $i=1$ or $h_i$ is immediately preceded by a vertical edge.
\end{corollary}
\begin{proof}
 By Proposition~\ref{prop:piecewise equivalence} and Proposition~\ref{prop:compatibility equivalence}, the grading $(\omega_{H'},\omega_{V'})\coloneqq \big((\varphi^*_{m-1})^{-1}\omega_V,\Omega_{m-1}^{-1}\omega_H)$ of $D'_{m-1}$ is piecewise compatible but not compatible.
 Let $h'_j\in H'_{m-1}$ denote the blocking edge of $\omega_{H'}$.
 Then since $m\ge5$, we have $|\rsh(h'_j;\omega_{H'})|=d_2-\ell$, where $\ell$ is the number of vertical edges immediately following $h'_j$.
 By Lemma~\ref{le:remote shadow cardinalities}, this implies there are $d_2-\ell$ horizontal edges of height $j-1$ in the remote shadow of $\omega_V$.
 But there are exactly $d_2-\ell$ horizontal edges of height $j-1$ inside $D_m$ by Lemma~\ref{le:Dyck path recursion}.
 Since $D(v_{u_{m-1,2}};\omega_V)=h_iv_{u_{m-1,2}}$, the edge $h_i$ lies furthest to the left among all horizontal edges in $\rsh(\omega_V)\cap h_iv_{u_{m-1,2}}$, this gives the result.
\end{proof}

We also obtain the following analogue of Proposition~\ref{prop:strong justification implication}.
\begin{corollary}
 \label{cor:support and remote shadow}
 Let $\omega:E_m\to\ZZ_{\ge0}$ be a piecewise compatible grading of $D_m$, $m\ge3$, which is not compatible.
 If $h_i\in H_m$ denotes the blocking edge for $\omega_H$, then $\supp(\omega_V)\cap h_iv_{u_{m-1,2}}=\rsh(\omega_H)\cap h_iv_{u_{m-1,2}}$.
\end{corollary}
\begin{proof}
 Since $\omega$ is not compatible, the grading $(\Omega_m(\omega_V),\varphi^*_m\omega_H)=:(\omega_{H'},\omega_{V'})$ of $D'_{m+1}$ is not compatible by Proposition~\ref{prop:compatibility equivalence}, but is piecewise compatible by Proposition~\ref{prop:piecewise equivalence}.
 By Theorem~\ref{th:blocking edge conditions} the grading $(\omega_{H'},\omega_{V'})$ satisfies the hypotheses of Proposition~\ref{prop:strong justification implication} and so $\supp(\omega_{H'})\cap h'_jv'_{u'_{m,2}}=\rsh(\omega_{V'})\cap h'_jv'_{u'_{m,2}}$, where $h'_j$ denotes the blocking edge of $\omega_{H'}$.

 By piecewise compatibility, we must have $\supp(\omega_V)\cap h_iv_{u_{m-1,2}}\subset\rsh(\omega_H)\cap h_iv_{u_{m-1,2}}$ since every vertical edge in $\supp(\omega_V)\cap h_iv_{u_{m-1,2}}$ is contained in the shadow of $\omega_H$.
 If there exists $v\in\rsh(\omega_H)\cap h_iv_{u_{m-1,2}}$ with $\omega_V(v)=0$, by Lemma~\ref{le:remote shadow cardinalities} there will be a horizontal edge $h'\in\rsh(\omega_{V'})\cap h'_jv'_{u'_{m,2}}$ with $\omega_{H'}(h')=0$, a contradiction.
 Therefore we must have $\supp(\omega_V)\cap h_iv_{u_{m-1,2}}=\rsh(\omega_H)\cap h_iv_{u_{m-1,2}}$.
\end{proof}

As a final consequence we show that the piecewise compatible gradings which are not compatible satisfy a certain upper bound property with respect to compatible gradings.
\begin{corollary}
 \label{cor:maximality}
 Suppose $\omega:E_m\to\ZZ_{\ge0}$ is a piecewise compatible grading of $D_m$, $m\ge3$, which is not compatible.
 Write $h_i$ for the blocking edge of $\omega_H$.
 Then the following hold:
 \begin{enumerate}[label=(\alph*)]
 \item\label{coro3.24_1} for any vertical grading $\chi_V\in\cC(\omega_H)$ and any edge $v\in(h_iv_{u_{m-1,2}})_V$, we have $\chi_V(v)\le\omega_V(v)$;
 \item\label{coro3.24_2} for any horizontal grading $\chi_H\in\cC(\omega_V)$ and any edge $h\in(h_iv_{u_{m-1,2}})_H$, we have $\chi_H(h)\le\omega_H(h)$.
 \end{enumerate}
\end{corollary}
\begin{proof}
 We begin by making a few basic observations which allow to deduce part~\ref{coro3.24_2} for $D_m$ from part~\ref{coro3.24_1} for $D'_{m-1}$.
 
 Consider a horizontal grading $\chi_H\in\cC(\omega_V)$ and suppose $\omega_H(h)<\chi_H(h)$ for some $h\in(h_iv_{u_{m-1,2}})_H$.
 This implies $\omega_H(h)<d_1$ since we only consider $d_1$-bounded horizontal gradings.
 By Theorem~\ref{th:blocking edge conditions}, we have $D(v_{u_{m-1,2}};\omega_V)=h_iv_{u_{m-1,2}}$ and so every edge of $h_iv_{u_{m-1,2}}$ is in the shadow of $\omega_V$.
 Thus we have 
 \[\supp(\chi_H)\cap h_iv_{u_{m-1,2}}\subset\rsh(\omega_V)\cap h_iv_{u_{m-1,2}}=\supp(\omega_H)\cap h_iv_{u_{m-1,2}},\]
 where the equality comes from Proposition~\ref{prop:strong justification implication}.
 By Theorem~\ref{th:blocking edge conditions}, $\omega_H$ is strongly left-justified at $h_i$ and so the only edge $h\in\supp(\omega_H)\cap h_iv_{u_{m-1,2}}$ which could satisfy $\omega_H(h)<d_1$ is $h=h_{i-1+d}$, where $d=|\supp(\omega_H)\cap h_iv_{u_{m-1,2}}|$.

 For $m=3$, we have $\supp(\omega_H)\cap h_iv_{u_{m-1,2}}=\{h_i\}$.
 Since the horizontal grading condition~\eqref{eq:hgc} of $\omega_H$ is not satisfied for the path $h_iv_{u_{m-1,2}}$, the inequality $\omega_H(h_i)<\chi_H(h_i)$ implies the horizontal grading condition of $\chi_H$ is also not satisfied for the path $h_iv_{u_{m-1,2}}$.
 In particular, $(\chi_H,\omega_V)$ is not compatible, a contradiction. 
 
 For $m\ge4$, consider the compatible grading $(\omega_{H'},\chi_{V'})\coloneqq \big((\varphi^*_{m-1})^{-1}\omega_V,\Omega_{m-1}^{-1}\chi_H\big)$ of $D'_{m-1}$ (see Proposition~\ref{prop:compatibility equivalence}) and the piecewise compatible grading $(\omega_{H'},\omega_{V'})\coloneqq \big((\varphi^*_{m-1})^{-1}\omega_V,\Omega_{m-1}^{-1}\omega_H\big)$ of $D'_{m-1}$ (see Proposition~\ref{prop:piecewise equivalence}).
 By the definition of $\Omega$, we have
 \[\chi_{V'}(\theta^{-1}h_{i-1+d})=\chi_H(h_{i-1+d})>\omega_H(h_{i-1+d})=\omega_{H'}(\theta^{-1}h_{i-1+d}).\]
 This contradicts part~\ref{coro3.24_1} applied to the grading $(\omega_{H'},\omega_{V'})$ of $D'_{m-1}$ and so there can be no grading $\chi_H$ as above.
 Thus part~\ref{coro3.24_2} holds for $m$ once we have established part~\ref{coro3.24_1} for $m-1$, $m\ge4$.

 To continue we suppose there exists a vertical grading $\chi_V\in\cC(\omega_H)$ such that $\chi_V(v)>\omega_V(v)$ for some $v\in(h_iv_{u_{m-1,2}})_V$.
 As above, this implies $0<\omega_V(v)<d_2$ and thus $v=v_{u_{m-1,2}-t+1}$, where $t=|\supp(\omega_V)\cap h_iv_{u_{m-1,2}}|$.
 In particular, we must have $d_2\ge2$ and by Corollary~\ref{cor:edge restrictions} the Dyck path $D_m$ has no consecutive vertical edges.
 
 Note that, by Proposition~\ref{prop:strong justification implication}, there are only two possibilities for the height of the edge $h_{i-1+d}$.
 Either $\hgt(h_{i-1+d})=u_{m-1,2}-t$ so that $v_{u_{m-1,2}-t+1}\in\rsh(h_{i-1+d};\omega_H)$ or $\hgt(h_{i-1+d})=u_{m-1,2}-t-1$ with $h_{i-1+d}$ immediately followed by a single vertical edge.
 In the latter case, $\omega_H(h_{i-1+d})>1$ also implies $v_{u_{m-1,2}-t+1}\in\rsh(h_{i-1+d};\omega_H)$.

 If $v_{u_{m-1,2}-t+1}\in\rsh(h_{i-1+d};\omega_H)$, the horizontal grading condition~\eqref{eq:hgc} is not satisfied for the path $h_{i-1+d}v_{u_{m-1,2}-t+1}$ and we have $D(v_{u_{m-1,2}-t+1};\omega_V)=h_{i+d}v_{u_{m-1,2}-t+1}$ by Proposition~\ref{prop:strong justification implication}.
 But then for $\chi_V$ as above, the vertical grading condition~\eqref{eq:vgc} is not satisfied for the path $h_{i-1+d}v_{u_{m-1,2}-t+1}$.
 In particular, this implies $(\omega_H,\chi_V)$ is not compatible, a contradiction.

 Thus we must have $\hgt(h_{i-1+d})=u_{m-1,2}-t-1$ with $h_{i-1+d}$ immediately followed by exactly one vertical edge and $\omega_H(h_{i-1+d})=1$.
 Then, since $\omega_H$ is strongly left-justified at $h_i$, we have $v_{u_{m-1,2}-t+1}\in\rsh(h_{i-2+d};\omega_H)$ and so the horizontal grading condition~\eqref{eq:hgc} is not satisfied for the path $h_{i-2+d}v_{u_{m-1,2}-t+1}$.
 If $h_{i-1+d}\in\rsh(v_{u_{m-1,2}-t+1};\omega_V)$, we must have $D(v_{u_{m-1,2}-t+1};\omega_V)=h_{i-1+d}v_{u_{m-1,2}-t+1}$.
 But then for $\chi_V$ as above, the vertical grading condition~\eqref{eq:vgc} is not satisfied for the path $h_{i-2+d}v_{u_{m-1,2}-t+1}$.
 In particular, this implies $(\omega_H,\chi_V)$ is not compatible, a contradiction.

 Thus the horizontal edge $h_{i-1+d}$ must lie beyond the shadow of $v_{u_{m-1,2}-t+1}$.
 By Proposition~\ref{prop:strong justification implication}, there can be no horizontal edges of height $u_{m-1,2}-t$ in the remote shadow of $\omega_V$ and so we must have $\omega_V(v_{u_{m-1,2}-t+1})=\ell$, where $\ell<d_2$ is the number of horizontal edges immediately preceding $v_{u_{m-1,2}-t+1}$. 
 For $m=3$, this can only occur for $t=1$, but $v_{u_{m-1,2}}$ is immediately preceded by $d_2-1$ horizontal edges inside $D_3$ and thus $\omega$ is compatible, a contradiction.

 So we must have $m\ge4$.
 Consider the piecewise compatible grading $(\omega_{H'},\omega_{V'})\coloneqq \big((\varphi^*_{m-1})^{-1}\omega_V,\Omega_{m-1}^{-1}\omega_H\big)$ of $D'_{m-1}$ (see Proposition~\ref{prop:piecewise equivalence}).
 Since $\omega_V$ is strongly right-justified and $\omega_V(v_{u_{m-1,2}-t+1})<d_2$, the last horizontal edge in $\supp(\omega_{H'})$ must be $h'_{u_{m-1,2}-t+1}$ with $\omega_{H'}(h'_{u_{m-1,2}-t+1})=d_2-\ell$, this being exactly the number of vertical edges immediately following $h'_{u_{m-1,2}-t+1}$ by Lemma~\ref{le:Dyck path recursion}.
 Moreover, by Lemma~\ref{le:remote shadow cardinalities}, the first vertical edge $v'$ in $\rsh(\omega_{H'})$ lies in the remote shadow of $h'_{u_{m-1,2}-t}$ and $\omega_{V'}(v')=1$.
 By Corollary~\ref{cor:support and remote shadow}, $v'$ cannot be immediately preceded by a vertical edge.
 But then there exists $\chi_{V'}$ with $\chi_{V'}(v')=2$ compatible with $\omega_{H'}$, a contradiction with part~\ref{coro3.24_2} for $D'_{m-1}$.
 
 This contradiction shows there can be no vertical $\chi_V$ as above and thus proves~\ref{coro3.24_1}.
\end{proof}


%%%%%%%%%%%%%%%%%%%%%%%%%%%%%%%%%%%%%%%%%%%%%%%%%%%%%%%%
\section{Proof of Main Theorem}
\label{sec:proof of main}

\noindent
We begin this section with a general statement about non-commutative weights associated to certain gradings of an arbitrary (\ie not necessarily maximal) Dyck path, here we make no boundedness assumptions on the gradings.

\begin{proposition}
 \label{prop:noncommutative collapse}
 Let $D$ be any Dyck path with edges $E=H\sqcup V$, where $H=\{h_1,\ldots,h_{a_1}\}$ with $a_1\ge1$ and $V=\{v_1,\ldots,v_{a_2}\}$ denote the sets of horizontal and vertical edges of $D$. 
 Write $E=\{1,2,\ldots,a_1+a_2\}$ for the edges of $D$ taken in the natural order.
 Let $\omega:E\to\ZZ_{\ge0}$ be any grading of $D$.
 Given $q_{i,j}\in\kk$ for $i\in\{1,2\}$ and $j\in\ZZ_{\ge0}$, define non-commutative weights
 \begin{equation}
 \label{eq:general edge weights}
 \wt_\omega(e)=\begin{cases}
 q_{1,\omega(e)}Y^{\omega(e)}X^{-1} & \text{if $e\in H$;}\\
 q_{2,\omega(e)}X^{\omega(e)+1}Y^{-1}X^{-1} & \text{if $e\in V$;}\\
 \end{cases}
 \end{equation}
 and let $Y_D(\omega)=\wt_\omega(1)\wt_\omega(2)\cdots\wt_\omega(a_1+a_2)$.
 Assume $\omega$ is compatible and satisfies the following:
 \begin{enumerate}
 \item\label{prop4.1_1} the local shadow path $D(h_1;\omega_H)=D$ with $f_{\omega_H}(D)=0$;
 \item\label{prop4.1_2} for any other vertical grading $\chi_V\in\cC(\omega_H)$ and any vertical edge $v_t\in V$ so that $\chi_V(v_s)=\omega_V(v_s)$ for $s<t$, we have $\chi_V(v_t)\le\omega_V(v_t)$.
 \end{enumerate}
 Then $Y_D(\omega)=pX^{-1}$, where $p=\prod_{i=1}^{a_1}q_{1,\omega(h_i)}\cdot\prod_{t=1}^{a_2}q_{2,\omega(v_t)}$.
\end{proposition}
\begin{proof}
 We first note that the coefficient $p$ is immediate from the definition of the non-commutative edge weights in equation~\eqref{eq:general edge weights}. 
 Thus we assume all $q_{i,j}=1$ for the remainder of the proof.

 We work by induction on $a_2$. 
 For $a_2=0$, assumption~\ref{prop4.1_1} implies $a_1=1$ and $\omega_H(h_1)=0$. 
 The claim follows in this case directly from the definition of the non-commutative edge weights in equation~\eqref{eq:general edge weights}.

 Suppose $a_2\ge1$ and consider $h_i\in\supp(\omega_H)$ with $i$ maximal.
 Let $v_r\in V$ denote the next vertical edge after $h_i$, \ie the path $h_iv_r$ consists of several consecutive horizontal edges, say $d$ of them, followed by a single vertical edge.
 By assumption~\ref{prop4.1_1}, we have $r<a_2$.
 By assumption~\ref{prop4.1_2}, we have $\omega_V(v_r)=d-1$ so that 
 \begin{equation}\label{dagger2}
 \tag{$\dagger$}
 Y_{h_iv_r}(\omega)=\big(Y^{\omega_H(h_i)}X^{-1}\big)\big(X^{-1}\big)^{d-1}\big(X^dY^{-1}X^{-1}\big)=Y^{\omega_H(h_i)-1}X^{-1}.
 \end{equation}
 Let $\widetilde D$ be the Dyck path obtained from $D$ by replacing the path $h_iv_r$ by a single horizontal edge.
 Write $\widetilde E=\widetilde H\sqcup\widetilde V$ for the edges of $\widetilde D$, where $\widetilde H=\{\tilde h_1,\ldots,\tilde h_{a_1-d+1}\}$ and $\widetilde V=\{\tilde v_1,\ldots,\tilde v_{a_2-1}\}$ denote the horizontal and vertical edges of $\widetilde D$.
 Define a grading $\widetilde\omega:\widetilde E\to\ZZ_{\ge0}$ by
 \[\widetilde\omega_H(\tilde h_j)=\begin{cases} \omega_H(h_j) & \text{if $j<i$;}\\ \omega_H(h_i)-1 & \text{if $j=i$;}\\ 0 & \text{if $j>i$;}\end{cases}
 \hspace{1in}
 \widetilde\omega_V(\tilde v_s)=\begin{cases} \omega_V(v_s) & \text{if $s<r$;}\\ \omega_V(v_{s+1}) & \text{if $s\ge r$.}\end{cases}\]
 It is not hard to see that $\widetilde\omega$ satisfies assumptions~\ref{prop4.1_1} and~\ref{prop4.1_2}, thus by induction we have $Y_{\widetilde D}(\widetilde\omega)=X^{-1}$.
 By~\eqref{dagger2}, we have $Y_D(\omega)=Y_{\widetilde D}(\widetilde\omega)$ and so $Y_D(\omega)=X^{-1}$ as desired.
\end{proof}

Now we turn to the proof of Theorem~\ref{th:combinatorial construction} and return to our standard boundedness assumptions on gradings.
\begin{lemma}
 \label{le:noncommutative collapse}
 Let $\omega:E_m\to\ZZ_{\ge0}$ be a piecewise compatible grading of $D_m$, $m\ge3$, which is not compatible.
 Denote by $h_i\in H_m$ the blocking edge of $\omega_H$. 
 Set 
 \[d=|\supp(\omega_H)\cap h_iv_{u_{m-1,2}}|\quad\text{ and }\quad t=|\supp(\omega_V)\cap h_iv_{u_{m-1,2}}|.\]
 Then for any $h\in(\overline{h}_ih_{i-1+d})_H$, we have $Y_{D(h;\omega_H)}(\omega)=pX^{-1}$, where $p=p_{1,\omega_H(h_{i-1+d})}p_{2,d_2-\omega_V(v_{u_{m-1,2}-t+1})}$.
\end{lemma}
\begin{proof}
 Since $h_i$ is blocking, no local shadow path $D(h;\omega_H)$ for $h\in(\overline{h}_ih_{i-1+d})_H$ contains $v_{u_{m-1,2}}$.
 Thus Lemma~\ref{le:compatibility check} shows $\omega|_{D(h;\omega_H)}$ is compatible.
 By definition of local shadow paths, $\omega|_{D(h;\omega_H)}$ satisfies condition~\ref{prop4.1_1} of Proposition~\ref{prop:noncommutative collapse}.
 Condition~\ref{prop4.1_2} follows directly from Corollary~\ref{cor:maximality}.
 The conclusion immediately follows since the only edges in $D(h;\omega_H)$ for $h\in(\overline{h}_ih_{i-1+d})_H$ whose non-commutative weights have nontrivial coefficients are $h_{i-1+d}$ and $v_{u_{m-1,2}-t+1}$.
\end{proof}

This leads to the following result which is key to our induction argument.
\begin{corollary}
 \label{cor:incompatible collapse}
 Let $\omega:E_m\to\ZZ_{\ge0}$ be a piecewise compatible grading of $D_m$, $m\ge3$, which is not compatible. 
 Write $h_i$ for the blocking edge of $\omega_H$ and assume $f_{\omega_H}(h_iv_{u_{m-1,2}})=0$. 
 Set 
 \[d=|\supp(\omega_H)\cap h_iv_{u_{m-1,2}}|\quad\text{ and }\quad t=|\supp(\omega_V)\cap h_iv_{u_{m-1,2}}|.\]
 Then $Y_{h_iv_{u_{m-1,2}}}(\omega)=pYXY^{-1}X^{-1}$, where $p=p_{1,\omega_H(h_{i-1+d})}p_{2,d_2-\omega_V(v_{u_{m-1,2}-t+1})}$.
\end{corollary}
\begin{proof}
 We distinguish two cases as in the proof of Theorem~\ref{th:blocking edge conditions}.
 First consider the case $\hgt(h_i)\ge u_{m-1,2}-d_1$ in which one of the following holds: $m=3$ with $d_1\ne1$, $m=4$ with $d_1,d_2\ne1$, or $m=5$ with $d_2=1$.
 In each of these cases $\supp(\omega_H)\cap h_iv_{u_{m-1,2}}=\{h_i\}$ and by assumption $\omega_H(h_i)=u_{m-1,2}-\hgt(h_i)$.
 We use the description of $\omega$ from the proof of Theorem~\ref{th:blocking edge conditions} in each case.

 For $m=3$ with $d_1\ne1$, set $\delta=1$ if $h_i$ is immediately followed by a vertical edge and $\delta=0$ otherwise. 
 Then we have
 \begin{align*}
 Y_{h_iv_{u_{2,2}}}(\omega)
 &=\big(p_{1,\omega_H(h_i)}Y^{\omega_H(h_i)} X^{-1}\big)\big(XY^{-1}X^{-1}\big)^\delta\big(X^{-1}\big)^{\omega_V(v_{u_{m-1,2}-t+1})}\\
 &\quad\times\big(p_{2,d_2-\omega_V(v_{u_{m-1,2}-t+1})}X^{\omega_V(v_{u_{m-1,2}-t+1})+1}Y^{-1}X^{-1}\big)\\
 &\quad\times\Big[\big(X^{-1}\big)^{d_2}\big(X^{d_2+1}Y^{-1}X^{-1}\big)\Big]^{\omega_H(h_i)-2-\delta}\big(X^{-1}\big)^{d_2-1}\big(X^{d_2+1}Y^{-1}X^{-1}\big)\\
 &=p\big(Y^{\omega_H(h_i)-1-\delta}X^{-1}\big)\Big[XY^{-\omega_H(h_i)+2+\delta}X^{-1}\Big]\big(X^2Y^{-1}X^{-1}\big)\\
 &=pYXY^{-1}X^{-1}.
 \end{align*}
 For $m=4$ with $d_1,d_2\ne1$ or $m=5$ with $d_2=1$, we have $p=1$ and so
 \begin{align*}
 Y_{h_iv_{u_{m-1,2}}}(\omega)
 &=\big(Y^{d_1}X^{-1}\big)\big(XY^{-1}X^{-1}\big)^{1+\delta_1}\Big[\big(X^{-1}\big)^{d_2}\big(X^{d_2+1}Y^{-1}X^{-1}\big)\Big]^{d_1-2-\delta_1}\\
 &\quad\times\big(X^{-1}\big)^{d_2-1}\big(X^{d_2+1}Y^{-1}X^{-1}\big)\\
 &=\big(Y^{d_1-1-\delta_1}X^{-1}\big)\Big[XY^{-d_1+2+\delta_1}X^{-1}\Big]\big(X^2Y^{-1}X^{-1}\big)\\
 &=YXY^{-1}X^{-1}.
 \end{align*}

 Now suppose $\hgt(h_i)<u_{m-1,2}-d_1$ so that $|\supp(\omega_H)\cap h_iv_{u_{m-1,2}}|>1$ and, by Theorem~\ref{th:blocking edge conditions}\ref{theo3.20_2_c}, $\hgt(h_{i+1})=\hgt(h_i)+\delta_1$.
 By Lemma~\ref{le:noncommutative collapse}, we have $Y_{D(h_{i+1};\omega_H)}(\omega)=pX^{-1}$.
 As in the proof of Theorem~\ref{th:blocking edge conditions}, we have $\omega_H(h_i)=d_1$ and $\omega_V(v_{u_{m-1,2}-\ell})=d_2$ for $0\le\ell\le d_1-\delta_1$.
 Combining these observations, we get
 \begin{align*}
 Y_{h_iv_{u+{m-1,2}}}(\omega) 
 &=\big(Y^{d_1}X^{-1}\big)\big(XY^{-1}X^{-1}\big)^{\delta_1}\big[pX^{-1}\big]\Big[\big(X^{-1}\big)^{d_2-1}\big(X^{d_2+1}Y^{-1}X^{-1}\big)\Big]\\
 &\quad\times\Big[\big(X^{-1}\big)^{d_2}\big(X^{d_2+1}Y^{-1}X^{-1}\big)\Big]^{d_1-2-\delta_1}\big(X^{-1}\big)^{d_2-1}\big(X^{d_2+1}Y^{-1}X^{-1}\big)\\
 &=p\big(Y^{d_1-1-\delta_1}X^{-1}\big)\Big[XY^{-d_1+2+\delta_1}X^{-1}\Big]\big(X^2Y^{-1}X^{-1}\big)\\
 &=pYXY^{-1}X^{-1}.\qedhere
 \end{align*}
\end{proof}

For $m\ge1$, we consider summands of $Y_{D_m}$ given as follows:
\begin{equation}
 \label{eq:Y decomposition}
 Y_{D_m}=\sum_{\omega_H:H_m\to[0,d_1]}Y_{D_m}(\omega_H),\qquad Y_{D_m}(\omega_H)\coloneqq \sum_{\omega_V\in\cC(\omega_H)}Y_{D_m}(\omega_H,\omega_V).
\end{equation}
Our goal will be to understand the action of $F_{P_0}$ on each of these summands.
The first step is given by the following factorization results which allow for an induction argument.
\begin{lemma}
 \label{le:piecewise compatible factorization 1}
 Let $\omega_H:H_m\to[0,d_1]$ be a horizontal grading of $D_m$, $m\ge3$.
 Write 
 \[Y_{D_m}^{pc}(\omega_H)=\sum\limits_{\omega:E_m\to\ZZ_{\ge0}}Y_{D_m}(\omega),\]
 where the sum ranges over piecewise compatible gradings $\omega$ of $D_m$ for which $\omega|_{H_m}=\omega_H$.
 Then there is the following factorization:
 \begin{multline}
 \label{eq:piecewise compatible factorization 1}
 Y_{D_m}^{pc}(\omega_H)\\
 =\Mk Y_{\mk D_{m\mk -\mk 1}}(\omega_{H,1})Y_{D_{m\mk -\mk 1}}(\omega_{H,2})\cdots Y_{D_{m\mk -\mk 1}}(\omega_{H,d_m\mk -\mk 1\mk -\mk \delta_m})pX^{|H_{m\mk -\mk 2\mk -\mk \delta_m}|} Y_{D_{m\mk -\mk 1}}(\omega_{H,d_m\mk -\mk \delta_m}),
 \end{multline}
 where 
 \[p=\begin{cases} p_{1,1}^{|V_{m-2-\delta_m}|-2|H_{m-2-\delta_m}|}p_{1,2}^{|H_{m-2-\delta_m}|-|V_{m-2-\delta_m}|} & \text{if $d_2=1$ and $m>3$;}\\ p_{1,1}^{-|V_{m-2-\delta_m}|} & \text{if $d_2>1$ or $m=3$.}\end{cases}\]
\end{lemma}
\begin{proof}
 By the assumptions on the horizontal grading $\omega_{H,d_m-\delta_m}$ of $D_{m-1}$ from Definition~\ref{def:agregate grading}, each term contributing to $Y_{D_{m-1}}(\omega_{H,d_m-\delta_m})$ begins with the monomial $p^{-1}X^{-|H_{m-2-\delta_m}|}$ associated to the initial $D_{m-2-\delta_m}$ subpath of $D_{m-1}$. 
 To see the coefficient $p^{-1}$, we observe the following:
 \begin{itemize}
 \item when $d_2>1$, there are $|V_{m-2-\delta_m}|$ horizontal edges of $D_{m-2-\delta_m}$ which are immediately followed by a single vertical edge and all other horizontal edges are not immediately followed by any vertical edges;
 \item when $d_2=1$, each horizontal edge is immediately followed by a vertical edge and so there are $|V_{m-2-\delta_m}|-|H_{m-2-\delta_m}|$ horizontal edges of $D_{m-2-\delta_m}$ which are immediately followed by exactly two vertical edges (see Corollary~\ref{cor:edge restrictions}) and the remaining $2|H_{m-2-\delta_m}|-|V_{m-2-\delta_m}|$ horizontal edges are immediately followed by a single vertical edge.
 \end{itemize}

 Using the notation of Definition~\ref{def:subpath edges}, for any grading $\omega:E_m\to\ZZ_{\ge0}$ there is the factorization
 \begin{multline*}
 Y_{D_m}(\omega)=Y_{h_{1,1}v_{u_{m-2,2},1}}(\omega)\cdots Y_{h_{1,d_m-1-\delta_m}v_{u_{m-2,2},d_m-1-\delta_m}}(\omega)\\
% \times 
 Y_{h_{u_{m-2-\delta_m,1}+1,d_m-\delta_m}v_{u_{m-2,2},d_m-\delta_m}}(\omega).
 \end{multline*}
 The result then immediately follows from the definition of piecewise compatible gradings in Definition~\ref{def:agregate grading}.
\end{proof}

Using Remark~\ref{rem:alternate restrictions} instead of Definition~\ref{def:agregate grading}, we obtain a similar factorization for piecewise compatible gradings of $D'_{m+1}$.
Below we use the notation $Y'_{D'_m}(\omega_{V'})\coloneqq \sum_{\omega_{H'}\in\cC(\omega_{V'})}Y'_{D'_m}(\omega_{H'},\omega_{V'})$ for a vertical grading $\omega_{V'}:V'_m\to[0,d_1]$.
Note that $d'_{m+1}=d_m$, $d'_m=d_{m-1}$, and so $\delta'_{m+1}=\delta_m$ when $m\ge3+\delta'_{m+1}$.
\begin{lemma}
 \label{le:piecewise compatible factorization 2}
 Let $\omega_{V'}:V'_{m+1}\to[0,d_1]$ be a vertical grading of $D'_{m+1}$ for $m\ge3+\delta'_{m+1}$.
 Write 
 \[{Y'}_{D'_{m+1}}^{pc}(\omega_{V'})=\sum\limits_{\omega':E'_{m+1}\to\ZZ_{\ge0}} Y'_{D'_{m+1}}(\omega'),\]
 where the sum ranges over piecewise compatible gradings $\omega'$ of $D'_{m+1}$ for which $\omega'|_{V'_{m+1}}=\omega_{V'}$.
 Then there is the following factorization:
 \begin{multline}
 \label{eq:piecewise compatible factorization 2}
 {Y'}_{D'_{m+1}}^{pc}(\omega_{V'})\\
 =\Mk Y'_{D'_m}\Mk (\omega_{V',1})Y'_{D'_m}\Mk (\omega_{V',2})\mk \cdots \mk Y'_{D'_m}(\omega_{V',d_m\mk -\mk 1\mk -\mk \delta_m})pXY^{|V'_{m\mk -\mk 1\mk -\mk \delta_m}|}X^{-1} Y'_{D'_m}(\omega_{V',d_m-\delta_m}),
 \end{multline}
 where 
 \[p=\begin{cases} p_{1,1}^{|H_{m-2-\delta_m}|-2|H_{m-2-\delta_m}|}p_{1,2}^{|H_{m-2-\delta_m}|-|V_{m-2-\delta_m}|} & \text{if $d_2=1$;}\\ p_{1,1}^{-|V_{m-2-\delta_m}|} & \text{if $d_2>1$.}\end{cases}\]
\end{lemma}
\begin{proof}
 By the assumptions on the vertical grading $\omega_{V',d_m-\delta_m}$ from Remark~\ref{rem:alternate restrictions}, each term contributing to $Y'_{D'_m}(\omega_{V',d_m\Mk -\Mk \delta_m})$ begins with the monomial $p^{\mk -\mk 1}\Mk XY^{-|V'_{m\mk -\mk 1\mk -\mk \delta_m}|}X^{\mk -\mk 1}$ associated to the initial $D'_{m-1-\delta_m}$ subpath of $D'_m$.
 The coefficient $p^{-1}$ here can be seen as follows. 
 Applying Lemma~\ref{le:Dyck path recursion}\ref{lemma2.3_1}, we see that the structure of $D'_{m-1-\delta_m}$ is determined by the structure of $D_{m-2-\delta_m}$ observed in the last part of the previous proof.
 More precisely, we have the following:
 \begin{itemize}
 \item when $d_2>1$, there are $|V_{m-2-\delta_m}|$ vertical edges of $D'_{m-1-\delta_m}$ which are immediately preceded by $d_1-1$ horizontal edges and all other vertical edges are immediately followed by $d_1$ horizontal edges;
 \item when $d_2=1$, there are $|V_{m-2-\delta_m}|-|H_{m-2-\delta_m}|$ vertical edges of $D'_{m-1-\delta_m}$ which are immediately preceded by $d_1-2$ horizontal edges and the remaining $2|H_{m-2-\delta_m}|-|V_{m-2-\delta_m}|$ vertical edges are immediately preceded by $d_1-1$ horizontal edges.
 \end{itemize}
 Then observe that in the computation of $Y'_{D'_m}(\omega_{V',d_m-\delta_m})$ the coefficients are given by $p'_{2,d_1-k}=p_{1,k}$ for $k=1,2$.
\end{proof}

The analogous factorization in the special case where $m=3$ and $\delta'_4=1$ is handled in the following result which is proven exactly as Lemma~\ref{le:piecewise compatible factorization 1}.
\begin{lemma}
 \label{le:piecewise compatible factorization 3}
 Suppose $d_2=1$. Let $\omega_{V'}:V'_4\to[0,d_1]$ be a vertical grading of $D'_4$.
 Write 
 \[{Y'}_{D'_4}^{pc}(\omega_{V'})=\sum\limits_{\omega':E'_4\to\ZZ_{\ge0}} Y'_{D'_4}(\omega'),\]
 where the sum ranges over piecewise compatible gradings $\omega'$ of $D'_4$ for which $\omega'|_{V'_4}=\omega_{V'}$.
 Then there is the following factorization:
 \begin{equation}
 \label{eq:piecewise compatible factorization 3}
 {Y'}_{D'_4}^{pc}(\omega_{V'})=Y'_{D'_3}(\omega_{V',1})Y'_{D'_3}(\omega_{V',2})\cdots Y'_{D'_3}(\omega_{V',d_1-2})X Y'_{D'_3}(\omega_{V',d_1-1}).
 \end{equation}
\end{lemma}

The factorizations above concerned sums over piecewise compatible gradings. Our goal is to understand sums over compatible gradings, however it will be easier to first focus on piecewise compatible gradings which are not compatible.
\begin{lemma}
 \label{le:incompatible factorization 1}
 Let $\omega_H:H_m\to[0,d_1]$ be a horizontal grading of $D_m$, $m\ge3$, for which there exists a vertical grading $\omega^*_V:V_m\to[0,d_2]$ of $D_m$ so that $(\omega_H,\omega^*_V)$ is piecewise compatible but not compatible.
 Write 
 \[Y_{D_m}^{nc}(\omega_H)=\sum\limits_{\omega:E_m\to\ZZ_{\ge0}}Y_{D_m}(\omega),\]
 where the sum ranges over piecewise compatible gradings $\omega$ of $D_m$ which are not compatible and satisfy $\omega|_{H_m}=\omega_H$.
 Let $h_i\in H_m$ denote the blocking edge of $\omega_H$ and set 
 \[d=|\supp(\omega_H)\cap h_iv_{u_{m-1,2}}|\quad\text{ and }\quad t=|\supp(\omega^*_V)\cap h_iv_{u_{m-1,2}}|.\]
 Let $s=\left\lceil\frac{i}{u_{m-1,1}}\right\rceil$ denote the index so that $h_i\in H_{m,s}$.
 Define a horizontal grading $\chi_H:H_{m-1}\to[0,d_1]$ of $D_{m-1}$ with $\chi_H(h)=\omega_{H,s}(h)$ for $h\in (h_{1,s}\overline{h}_i)_H$ and $\chi_H(h)=\ell$ for $h\in (h_iv_{u_{m-2,2},s})_H$ if $h$ is immediately followed by exactly $\ell$ vertical edges in this copy of $D_{m-1}$.
 Then there is the following factorization:
 \begin{multline}
 \label{eq:incompatible factorization 1}
 Y_{D_m}^{nc}(\omega_H)\\
 =Y_{D_{m-1}}(\omega_{H,1})\cdots Y_{D_{m-1}}(\omega_{H,s-1})Y_{D_{m-1}}(\chi_H)p_1X^{|h_iv_{u_{m-2,2},s}|_H} p_2YXY^{-1}X^{-1},
 \end{multline}
 where $p_2=p_{1,\omega_H(h_{i-1+d})}p_{2,d_2-\omega^*_V(v_{u_{m-1,2}-t+1})}$ and
 \[p_1=\begin{cases} p_{1,1}^{|h_iv_{u_{m-2,2},s}|_V-2|h_iv_{u_{m-2,2},s}|_H}p_{1,2}^{|h_iv_{u_{m-2,2},s}|_H-|h_iv_{u_{m-2,2},s}|_V} & \text{if $d_2=1$;}\\ 
 p_{1,1}^{-|h_iv_{u_{m-2,2},s}|_V} & \text{if $d_2>1$.}\end{cases}\]
\end{lemma}
\begin{proof}
 Using the notation of Definition~\ref{def:subpath edges}, for any grading $\omega:E_m\to\ZZ_{\ge0}$ there is the factorization
 \[
 Y_{D_m}(\omega)=Y_{h_{1,1}v_{u_{m-2,2},1}}(\omega)\cdots Y_{h_{1,s-1}v_{u_{m-2,2},s-1}}(\omega)Y_{h_{1,s}\overline{h}_i}(\omega)Y_{h_iv_{u_{m-1,2}}}(\omega).
 \]
 By definition of $\chi_H$, every term of $Y_{D_{m-1}}(\chi_H)$ ends with the monomial 
 \[
 p_1^{-1}X^{-|h_iv_{u_{m-2,2},s}|_H}.
 \]
 Theorem~\ref{th:blocking edge conditions} shows that any piecewise compatible grading $\omega:E_m\to\ZZ_{\ge0}$ of $D_m$ which is not compatible agrees with $(\omega_H,\omega^*_V)$ on the path $h_iv_{u_{m-1,2}}$.
 The result then immediately follows from Corollary~\ref{cor:incompatible collapse} and the definition of piecewise compatible gradings in Definition~\ref{def:agregate grading}.
\end{proof}

The next result gives an analogous factorization for sums over piecewise compatible gradings of $D'_{m+1}$ which are not compatible.
\begin{lemma}
 \label{le:incompatible factorization 2}
 Let $\omega_{V'}:V'_{m+1}\to[0,d_1]$ be a vertical grading of $D'_{m+1}$, $m\ge3$, for which there exists a horizontal grading $\omega^*_{H'}:H'_{m+1}\to[0,d_2]$ of $D'_{m+1}$ so that $(\omega^*_{H'},\omega_{V'})$ is piecewise compatible but not compatible.
 Write 
 \[{Y'}_{D'_{m+1}}^{nc}(\omega_{V'})=\sum\limits_{\omega:E'_{m+1}\to\ZZ_{\ge0}}Y'_{D'_{m+1}}(\omega),\]
 where the sum ranges over piecewise compatible gradings $\omega$ of $D'_{m+1}$ which are not compatible and satisfy $\omega|_{V'_{m+1}}=\omega_{V'}$.
 Let $h'_j\in H'_{m+1}$ denote the blocking edge of $\omega^*_{H'}$, where $\hgt(h'_j)=i-1$, and set 
 \[d=|\supp(\omega_{V'})\cap h'_jv'_{u'_{m,2}}|\quad\text{ and }\quad t=|\supp(\omega^*_H)\cap h'_jv'_{u'_{m,2}}|.\]
 Let $s=\left\lceil\frac{i}{u_{m-1,1}}\right\rceil$ denote the index so that $h'_j\in H'_{m+1,s}$.
 Define a horizontal grading $\chi_{V'}:V'_m\to[0,d_1]$ of $D'_m$ with $\chi_{V'}(v')=\omega_{V',s}(v')$ for $v'\in (h'_{1,s}\overline{h}'_j)_V$ and $\chi_{V'}(v')=\ell$ for $v'\in (h'_jv'_{u'_{m-1,2},s})_V$ if $v'$ is immediately preceded by exactly $\ell$ horizontal edges in this copy of $D'_m$.
 Then there is the following factorization:
 \begin{multline}
 \label{eq:incompatible factorization 2}
 {Y'}_{D'_{m+1}}^{nc}(\omega_{V'})\\
 =\Mk Y'_{D'_m}(\omega_{V',1})\cdots Y'_{D'_m}(\omega_{V',s-1})Y'_{D'_m}(\chi_{V'})p_1XY^{|h'_jv_{u'_{m-1,2},s}|_V}X^{-1} p_2YXY^{-1}X^{-1},
 \end{multline}
 where $p_2=p_{1,d_1-\omega_V(v'_{u'_{m,2}-d+1})}p_{2,d_2-\omega^*_H(h'_{j-1+t})}$ and
 \[p_1=\begin{cases} p_{1,1}^{|h_iv_{u_{m-2,2},s}|_V-2|h_iv_{u_{m-2,2},s}|_H}p_{1,2}^{|h_iv_{u_{m-2,2},s}|_H-|h_iv_{u_{m-2,2},s}|_V} & \text{if $d_2=1$;}\\ 
 p_{1,1}^{-|h_iv_{u_{m-2,2},s}|_V} & \text{if $d_2>1$;}\end{cases}\]
 with $h_iv_{u_{m-2,2,s}}$ being the subpath in the $s$-th copy of $D_{m-1}$ inside $D_m$.
\end{lemma}
\begin{proof}
 By definition of $\chi_{V'}$, every term of $Y'_{D'_m}(\chi_{V'})$ ends with the monomial $p_1^{-1}XY^{-|h'_jv'_{u'_{m-1,2},s}|_V}X^{-1}$.
 To see the coefficient $p_1^{-1}$, note that by Lemma~\ref{le:Dyck path recursion} the structure of $D'_m$ is determined by the structure of $D_{m-1}$.

 Theorem~\ref{th:blocking edge conditions} shows that any piecewise compatible grading $\omega:E'_{m+1}\to\ZZ_{\ge0}$ of $D'_{m+1}$ which is not compatible agrees with $(\omega^*_{H'},\omega_{V'})$ on the path $h'_jv'_{u'_{m,2}}$.
 The result then immediately follows from Corollary~\ref{cor:incompatible collapse} and the definition of piecewise compatible gradings in Definition~\ref{def:agregate grading}.
\end{proof}


